\input{preamble}

% Bien, commençons ici.
%
\begin{document}

\title{Champs de modules}

\maketitle

\phantomsection
\label{section-phantom}

\tableofcontents




\section{Introduction}
\label{section-introduction}

\noindent
Dans ce chapitre, nous vérifions des propriétés élémentaires d'espaces
et de champs de modules tels que
$\mathit{Hom}$, $\mathit{Isom}$, $\Cohstack_{X/B}$,
$\Quotfunctor_{\mathcal{F}/X/B}$, $\Hilbfunctor_{X/B}$,
$\Picardstack_{X/B}$, $\Picardfunctor_{X/B}$, $\mathit{Mor}_B(Z, X)$,
$\Spacesstack'_{fp, flat, proper}$, $\Polarizedstack$ et
$\Complexesstack_{X/B}$.
Nous avons déjà montré, sous des hypothèses convenables, qu'il s'agit
d'espaces algébriques ou de champs algébriques ; voir Quotients, sections
\ref{quot-section-hom},
\ref{quot-section-isom},
\ref{quot-section-stack-coherent-sheaves},
\ref{quot-section-not-flat},
\ref{quot-section-quot},
\ref{quot-section-hilb},
\ref{quot-section-picard-stack},
\ref{quot-section-picard-functor},
\ref{quot-section-relative-morphisms},
\ref{quot-section-stack-of-spaces},
\ref{quot-section-polarized} et
\ref{quot-section-moduli-complexes}.
Le champ des courbes, noté $\textit{Curves}$ et introduit dans
Quotients, section \ref{quot-section-curves}, est étudié au chapitre
sur les modules de courbes ; voir Modules des courbes, section
\ref{moduli-curves-section-stack-curves}.

\medskip\noindent
En un certain sens, ce chapitre suit les traces des exposés de
Grothendieck \cite{Gr-I},
\cite{Gr-II},
\cite{Gr-III},
\cite{Gr-IV},
\cite{Gr-V} et
\cite{Gr-VI}.






\section{Conventions et abus de langage}
\label{section-conventions}

\noindent
Nous continuons d'employer les conventions et les abus de langage
introduits dans Propriétés des champs, section
\ref{stacks-properties-section-conventions}.
Sauf mention contraire, notre schéma de base sera $\Spec(\mathbf{Z})$.






\section{Propriétés de Hom et Isom}
\label{section-hom-isom}

\noindent
Soit $f : X \to B$ un morphisme d'espaces algébriques de présentation
finie. Supposons que $\mathcal{F}$ et $\mathcal{G}$ soient des
$\mathcal{O}_X$-modules quasi-cohérents.
Si $\mathcal{G}$ est de présentation finie, plat sur $B$ et à support
propre sur $B$, alors le foncteur
$\mathit{Hom}(\mathcal{F}, \mathcal{G})$ défini par
$$
T/B \longmapsto \Hom_{\mathcal{O}_{X_T}}(\mathcal{F}_T, \mathcal{G}_T)
$$
est un espace algébrique affine sur $B$. Si $\mathcal{F}$ est de
présentation finie, alors
$\mathit{Hom}(\mathcal{F}, \mathcal{G}) \to B$
est de présentation finie. Voir Quotients, proposition
\ref{quot-proposition-hom}.

\medskip\noindent
Si $\mathcal{F}$ et $\mathcal{G}$ sont tous deux de présentation finie,
plats sur $B$ et à support propre sur $B$, alors le sous-foncteur
$$
\mathit{Isom}(\mathcal{F}, \mathcal{G}) \subset
\mathit{Hom}(\mathcal{F}, \mathcal{G})
$$
est un espace algébrique affine et de présentation finie sur $B$.
Voir Quotients, proposition \ref{quot-proposition-isom}.




\section{Propriétés du champ des faisceaux cohérents}
\label{section-stack-coherent-sheaves}

\noindent
Soit $f : X \to B$ un morphisme d'espaces algébriques séparé et de
présentation finie. Alors le champ $\Cohstack_{X/B}$ qui paramètre les
familles plates de modules cohérents à support propre est algébrique.
Voir Quotients, théorème
\ref{quot-theorem-coherent-algebraic-general}.

\begin{lemma}
\label{lemma-coherent-diagonal-affine-fp}
La diagonale de $\Cohstack_{X/B}$ sur $B$ est affine et de présentation
finie.
\end{lemma}

\begin{proof}
La représentabilité de la diagonale par des espaces algébriques a été
établie dans Quotients, lemme \ref{quot-lemma-coherent-diagonal}.
Il ressort de la démonstration qu'il nous faut montrer que
$\mathit{Isom}(\mathcal{F}, \mathcal{G}) \to T$
est affine et de présentation finie pour une paire de
$\mathcal{O}_{X_T}$-modules de présentation finie
$\mathcal{F}$, $\mathcal{G}$, plats sur $T$ et à support propre sur $T$.
C'est précisément ce qui a été étudié à la section
\ref{section-hom-isom}.
\end{proof}

\begin{lemma}
\label{lemma-coherent-qs-lfp}
Le morphisme $\Cohstack_{X/B} \to B$ est quasi-séparé et localement de
présentation finie.
\end{lemma}

\begin{proof}
Pour vérifier que $\Cohstack_{X/B} \to B$ est quasi-séparé, il faut
montrer que sa diagonale est quasi-compacte et quasi-séparée.
Cela résulte immédiatement du lemme
\ref{lemma-coherent-diagonal-affine-fp}.
Pour prouver que $\Cohstack_{X/B} \to B$ est localement de présentation
finie, il faut montrer que $\Cohstack_{X/B} \to B$ préserve les limites ;
voir Limites de champs, proposition
\ref{stacks-limits-proposition-characterize-locally-finite-presentation}.
Cela découle de Quotients, lemme \ref{quot-lemma-coherent-limits}
(un petit détail est omis).
\end{proof}

\begin{lemma}
\label{lemma-coherent-existence-part}
Supposons que $X \to B$ soit propre ainsi que de présentation finie.
Alors $\Cohstack_{X/B} \to B$ satisfait la partie d'existence du critère
valuatif (Morphismes de champs, définition
\ref{stacks-morphisms-definition-existence}).
\end{lemma}

\begin{proof}
Après changement de base, on se ramène immédiatement au problème suivant :
étant donnés un anneau de valuation $R$ de corps des fractions $K$, un espace
algébrique $X$ propre sur $R$ et un $\mathcal{O}_{X_K}$-module cohérent
$\mathcal{F}_K$, montrer qu'il existe un $\mathcal{O}_X$-module
$\mathcal{F}$ de présentation finie, plat sur $R$, dont la fibre générique
est $\mathcal{F}_K$.
Remarquons que, d'après Platitude sur les espaces, théorème
\ref{spaces-flat-theorem-finite-type-flat}, tout $\mathcal{O}_X$-module
quasi-cohérent $\mathcal{F}$ de type fini et plat sur $R$ est de
présentation finie.
Notons $j : X_K \to X$ l'immersion de la fibre générique.
Puisqu'il est obtenu par changement de base à partir du morphisme affine
$\Spec(K) \to \Spec(R)$, le morphisme $j$ est affine. Ainsi,
$j_*\mathcal{F}_K$ est quasi-cohérent. Écrivons
$$
j_*\mathcal{F}_K = \colim \mathcal{F}_i
$$
comme limite inductive filtrante de ses sous-$\mathcal{O}_X$-modules
quasi-cohérents de type fini ; voir Limites d'espaces, lemme
\ref{spaces-limits-lemma-directed-colimit-finite-type}.
Comme $j_*\mathcal{F}_K$ est un faisceau de $K$-espaces vectoriels sur
$X$, il est plat sur $\Spec(R)$. Chaque $\mathcal{F}_i$ est donc plat sur
$R$, car, sur un anneau de valuation, la platitude équivaut à l'absence de
torsion (Compléments d'algèbre, lemme
\ref{more-algebra-lemma-valuation-ring-torsion-free-flat}) et l'absence de
torsion passe aux sous-modules.
Enfin, nous devons montrer que l'application
$j^*\mathcal{F}_i \to \mathcal{F}_K$
est un isomorphisme pour un certain $i$.
Puisque $j^*j_*\mathcal{F}_K = \mathcal{F}_K$ (un petit détail est omis)
et que $j^*$ est exact, on voit que
$j^*\mathcal{F}_i \to \mathcal{F}_K$ est injective pour tout $i$.
Comme $j^*$ commute aux limites inductives, on a
$\mathcal{F}_K = j^*j_*\mathcal{F}_K = \colim j^*\mathcal{F}_i$.
Puisque $\mathcal{F}_K$ est cohérent (c'est-à-dire de présentation finie),
il existe un $i$ tel que $j^*\mathcal{F}_i$ contienne tous les générateurs,
en nombre fini, sur un recouvrement affine \'etale de $X$.
On obtient ainsi la surjectivité de
$j^*\mathcal{F}_i \to \mathcal{F}_K$ pour $i$ assez grand.
\end{proof}

\begin{lemma}
\label{lemma-coherent-functorial}
Soit $B$ un espace algébrique. Soit $\pi : X \to Y$ un morphisme
quasi-fini d'espaces algébriques séparés et de présentation finie sur $B$.
Alors $\pi_*$ induit un morphisme
$\Cohstack_{X/B} \to \Cohstack_{Y/B}$.
\end{lemma}

\begin{proof}
Soit $(T \to B, \mathcal{F})$ un objet de $\Cohstack_{X/B}$.
Nous affirmons que
\begin{enumerate}
\item[(a)] $(T \to B, \pi_{T, *}\mathcal{F})$ est un objet de
$\Cohstack_{Y/B}$, et que
\item[(b)] pour $T' \to T$, on a
$\pi_{T', *}(X_{T'} \to X_T)^*\mathcal{F} =
(Y_{T'} \to Y_T)^*\pi_{T, *}\mathcal{F}$.
\end{enumerate}
La partie (b) garantit que cette construction définit bien le foncteur
$\Cohstack_{X/B} \to \Cohstack_{Y/B}$ recherché.

\medskip\noindent
Soit $i : Z \to X_T$ le sous-espace fermé défini par le zéro-ième idéal
de Fitting de $\mathcal{F}$ (Diviseurs sur les espaces, section
\ref{spaces-divisors-section-fitting-ideals}).
Alors $Z \to B$ est propre par hypothèse (voir Catégories dérivées des
espaces, section \ref{spaces-perfect-section-proper-over-base}).
D'autre part, $i$ est de présentation finie (Diviseurs sur les espaces,
lemme \ref{spaces-divisors-lemma-fitting-ideal-of-finitely-presented}, et
Morphismes d'espaces, lemme
\ref{spaces-morphisms-lemma-closed-immersion-finite-presentation}).
Il existe un $\mathcal{O}_Z$-module quasi-cohérent $\mathcal{G}$ de type
fini tel que $i_*\mathcal{G} = \mathcal{F}$ (Diviseurs sur les espaces,
lemme \ref{spaces-divisors-lemma-on-subscheme-cut-out-by-Fit-0}).
En fait, $\mathcal{G}$ est de présentation finie comme
$\mathcal{O}_Z$-module, d'après Descente sur les espaces, lemme
\ref{spaces-descent-lemma-finite-finitely-presented-module}.
Remarquons que $\mathcal{G}$ est plat sur $B$, par exemple parce que les
germes de $\mathcal{G}$ et de $\mathcal{F}$ coïncident (Morphismes
d'espaces, lemme \ref{spaces-morphisms-lemma-stalk-push-closed}).
Remarquons que $\pi_T \circ i : Z \to Y_T$ est quasi-fini comme composée
de morphismes quasi-finis et que
$\pi_{T, *}\mathcal{F} = (\pi_T \circ i)_*\mathcal{G})$.
Puisque $i$ est affine, la formation de $i_*$ commute au changement de base
(Cohomologie des espaces, lemme
\ref{spaces-cohomology-lemma-affine-base-change}).
Nous pouvons donc remplacer $B$ par $T$, $X$ par $Z$,
$\mathcal{F}$ par $\mathcal{G}$ et $Y$ par $Y_T$ pour nous ramener au cas
étudié dans le paragraphe suivant.

\medskip\noindent
Supposons que $X \to B$ soit propre. Alors $\pi$ est propre d'après
Morphismes d'espaces, lemme
\ref{spaces-morphisms-lemma-universally-closed-permanence}, donc fini
d'après Compléments sur les morphismes d'espaces, lemme
\ref{spaces-more-morphisms-lemma-characterize-finite}.
Comme un morphisme fini est affine, la partie (b) résulte de Cohomologie
des espaces, lemme \ref{spaces-cohomology-lemma-affine-base-change}.
D'autre part, $\pi$ est de présentation finie d'après Morphismes d'espaces,
lemme \ref{spaces-morphisms-lemma-finite-presentation-permanence}.
Ainsi, $\pi_{T, *}\mathcal{F}$ est de présentation finie d'après Descente
sur les espaces, lemme
\ref{spaces-descent-lemma-finite-finitely-presented-module}.
Enfin, $\pi_{T, *}\mathcal{F} $ est plat sur $B$, comme on le voit par
exemple sur les germes en utilisant Cohomologie des espaces, lemme
\ref{spaces-cohomology-lemma-stalk-push-finite}.
\end{proof}

\begin{lemma}
\label{lemma-coherent-open}
Soit $B$ un espace algébrique. Soit $\pi : X \to Y$ une immersion ouverte
d'espaces algébriques séparés et de présentation finie sur $B$.
Alors le morphisme $\Cohstack_{X/B} \to \Cohstack_{Y/B}$ du lemme
\ref{lemma-coherent-functorial} est une immersion ouverte.
\end{lemma}

\begin{proof}
Omis. Indication : si $\mathcal{F}$ est un objet de $\Cohstack_{Y/B}$
sur $T$ et si, pour $t \in T$, on a
$\text{Supp}(\mathcal{F}_t) \subset |X_t|$, alors il en va de même pour
$t' \in T$ dans un voisinage de $t$.
\end{proof}

\begin{lemma}
\label{lemma-coherent-closed}
Soit $B$ un espace algébrique. Soit $\pi : X \to Y$ une immersion fermée
d'espaces algébriques séparés et de présentation finie sur $B$.
Alors le morphisme $\Cohstack_{X/B} \to \Cohstack_{Y/B}$ du lemme
\ref{lemma-coherent-functorial} est une immersion fermée.
\end{lemma}

\begin{proof}
Soit $\mathcal{I} \subset \mathcal{O}_Y$ le faisceau d'idéaux qui définit
$X$ comme sous-espace fermé de $Y$. Rappelons que $\pi_*$ induit une
équivalence entre la catégorie des $\mathcal{O}_X$-modules quasi-cohérents
et celle des $\mathcal{O}_Y$-modules quasi-cohérents annulés par
$\mathcal{I}$ ; voir Morphismes d'espaces, lemme
\ref{spaces-morphisms-lemma-i-star-equivalence}.
Il en va de même, mutatis mutandis, après changement de base par
$T \to B$, en remplaçant $\mathcal{I}$ par le faisceau d'idéaux
$\mathcal{I}_T = \Im((Y_T \to Y)^*\mathcal{I} \to \mathcal{O}_{Y_T})$.
L'analyse de la démonstration du lemme
\ref{lemma-coherent-functorial} montre que l'image essentielle de
$\Cohstack_{X/B} \to \Cohstack_{Y/B}$
est exactement formée des objets $\xi = (T \to B, \mathcal{F})$ pour
lesquels $\mathcal{F}$ est annulé par $\mathcal{I}_T$.
Autrement dit, $\xi$ appartient à l'image essentielle si et seulement si
l'application de multiplication
$$
\mathcal{F} \otimes_{\mathcal{O}_{Y_T}} (Y_T \to Y)^*\mathcal{I}
\longrightarrow
\mathcal{F}
$$
est nulle, de même qu'après tout changement de base ultérieur
$T' \to T$.
Notons que
$$
(Y_{T'} \to Y_T)^*(
\mathcal{F} \otimes_{\mathcal{O}_{Y_T}} (Y_T \to Y)^*\mathcal{I}) =
(Y_{T'} \to Y_T)^*\mathcal{F} \otimes_{\mathcal{O}_{Y_{T'}}}
(Y_{T'} \to Y)^*\mathcal{I})
$$
La nullité de l'application de multiplication sur $T'$ est donc
représentable par un sous-espace fermé de $T$, d'après Platitude sur les
espaces, lemme \ref{spaces-flat-lemma-F-zero-closed-proper}.
\end{proof}

\begin{situation}[Invariants numériques]
\label{situation-numerical}
Soit $f : X \to B$ comme dans l'introduction de cette section. Soit $I$
un ensemble et, pour $i \in I$, soit $E_i \in D(\mathcal{O}_X)$ parfait.
Étant donné un objet $(T \to B, \mathcal{F})$ de $\Cohstack_{X/B}$,
notons $E_{i, T}$ l'image inverse dérivée de $E_i$ sur $X_T$.
L'objet
$$
K_i = Rf_{T, *}(E_{i, T} \otimes_{\mathcal{O}_{X_T}}^\mathbf{L} \mathcal{F})
$$
de $D(\mathcal{O}_T)$ est parfait et sa formation commute au changement de
base ; voir Catégories dérivées des espaces, lemme
\ref{spaces-perfect-lemma-base-change-tensor-perfect}.
La fonction
$$
\chi_i : |T| \longrightarrow \mathbf{Z},\quad
\chi_i(t) =
\chi(X_t, E_{i, t} \otimes_{\mathcal{O}_{X_t}}^\mathbf{L} \mathcal{F}_t) =
\chi(K_i \otimes_{\mathcal{O}_T}^\mathbf{L} \kappa(t))
$$
est donc localement constante d'après Catégories dérivées des espaces,
lemme \ref{spaces-perfect-lemma-chi-locally-constant}.
Soit $P : I \to \mathbf{Z}$ une application. Considérons le sous-champ
$$
\Cohstack^P_{X/B} \subset \Cohstack_{X/B}
$$
formé des familles plates de faisceaux cohérents à support propre dont les
invariants numériques coïncident avec $P$. Plus précisément, un objet
$(T \to B, \mathcal{F})$ de $\Cohstack_{X/B}$ appartient à
$\Cohstack^P_{X/B}$ si et seulement si $\chi_i(t) = P(i)$ pour tout
$i \in I$ et tout $t \in T$.
\end{situation}

\begin{lemma}
\label{lemma-open-P}
Dans la situation \ref{situation-numerical}, le champ
$\Cohstack^P_{X/B}$ est algébrique et
$$
\Cohstack^P_{X/B} \longrightarrow \Cohstack_{X/B}
$$
est une immersion fermée plate. Si $I$ est fini ou si $B$ est localement
noethérien, alors $\Cohstack^P_{X/B}$ est un sous-champ ouvert et fermé de
$\Cohstack_{X/B}$.
\end{lemma}

\begin{proof}
C'est immédiat si $I$ est fini, puisque les fonctions
$t \mapsto \chi_i(t)$ sont localement constantes. Si $I$ est infini,
écrivons
$$
I = \bigcup\nolimits_{I' \subset I\text{ fini}} I'
$$
et notons $P' = P|_{I'}$. On a alors
$$
\Cohstack^P_{X/B} = \bigcap\nolimits_{I' \subset I\text{ fini}}
\Cohstack^{P'}_{X/B}
$$
Ainsi, $\Cohstack^P_{X/B}$ est toujours un champ algébrique et le morphisme
$\Cohstack^P_{X/B} \subset \Cohstack_{X/B}$ est toujours une immersion
fermée plate, mais il peut cesser d'être un sous-champ ouvert. (Nous
laissons au lecteur le soin de construire des exemples.) Toutefois, si
$B$ est localement noethérien, alors $\Cohstack_{X/B}$ l'est aussi, d'après
le lemme \ref{lemma-coherent-qs-lfp} et Morphismes de champs, lemme
\ref{stacks-morphisms-lemma-locally-finite-type-locally-noetherian}.
Par conséquent, si $U \to \Cohstack_{X/B}$ est un morphisme lisse
surjectif, où $U$ est un schéma localement noethérien, les images inverses
des sous-champs ouverts et fermés $\Cohstack^{P'}_{X/B}$ ont dans $U$ une
intersection ouverte (car les composantes connexes des espaces topologiques
localement noethériens sont ouvertes).
D'où le résultat dans ce cas.
\end{proof}

\begin{lemma}
\label{lemma-finite-list-perfect-objects}
Soit $f : X \to B$ comme dans l'introduction de cette section.
Soient $E_1, \ldots, E_r \in D(\mathcal{O}_X)$ parfaits.
Soit $I = \mathbf{Z}^{\oplus r}$ et considérons l'application
$$
I \longrightarrow D(\mathcal{O}_X),\quad
(n_1, \ldots, n_r) \longmapsto
E_1^{\otimes n_1}
\otimes \ldots \otimes
E_r^{\otimes n_r}
$$
Soit $P : I \to \mathbf{Z}$ une application. Alors
$\Cohstack^P_{X/B} \subset \Cohstack_{X/B}$,
tel qu'il est défini dans la situation \ref{situation-numerical}, est un
sous-champ ouvert et fermé.
\end{lemma}

\begin{proof}
Nous pouvons travailler localement pour la topologie \'etale sur $B$ et donc
supposer $B$ affine. Dans ce cas, nous pouvons effectuer une réduction
noethérienne absolue ; nous suggérons au lecteur de sauter la démonstration.
Écrivons en effet $B = \Spec(\Lambda)$.
Écrivons $\Lambda = \colim \Lambda_i$ comme limite inductive filtrante,
chaque $\Lambda_i$ étant de type fini sur $\mathbf{Z}$. Pour un certain
$i$, on peut trouver un morphisme d'espaces algébriques
$X_i \to \Spec(\Lambda_i)$ séparé et de présentation finie, dont le
changement de base à $\Lambda$ est $X$. Voir Limites d'espaces, lemmes
\ref{spaces-limits-lemma-descend-finite-presentation} et
\ref{spaces-limits-lemma-descend-separated-morphism}.
Après avoir augmenté $i$, on peut alors supposer qu'il existe des objets
parfaits $E_{1, i}, \ldots, E_{r, i}$ dans $D(\mathcal{O}_{X_i})$ dont
les images inverses dérivées sur $X$ sont respectivement isomorphes à
$E_1, \ldots, E_r$ ; voir Catégories dérivées des espaces, lemme
\ref{spaces-perfect-lemma-perfect-on-limit}.
On a manifestement un carré cartésien
$$
\xymatrix{
\Cohstack^P_{X/B} \ar[r] \ar[d] &
\Cohstack_{X/B} \ar[d] \\
\Cohstack^P_{X_i/\Spec(\Lambda_i)} \ar[r] &
\Cohstack_{X_i/\Spec(\Lambda_i)}
}
$$
et l'on peut donc appliquer le lemme \ref{lemma-open-P} pour achever la
démonstration.
\end{proof}

\begin{example}[Faisceaux cohérents à polynôme de Hilbert fixé]
\label{example-hilbert-polynomial}
Soit $f : X \to B$ comme dans l'introduction de cette section.
Soit $\mathcal{L}$ un $\mathcal{O}_X$-module inversible.
Soit $P : \mathbf{Z} \to \mathbf{Z}$ un polynôme numérique.
On peut alors considérer le sous-champ algébrique ouvert et fermé
$$
\Cohstack^P_{X/B} =
\Cohstack^{P, \mathcal{L}}_{X/B}
\subset \Cohstack_{X/B}
$$
formé des familles plates de faisceaux cohérents à support propre dont les
invariants numériques coïncident avec $P$ : un objet
$(T \to B, \mathcal{F})$ de $\Cohstack_{X/B}$ appartient à
$\Cohstack^P_{X/B}$ si et seulement si
$$
P(n) =
\chi(X_t, \mathcal{F}_t \otimes_{\mathcal{O}_{X_t}} \mathcal{L}_t^{\otimes n})
$$
pour tout $n \in \mathbf{Z}$ et tout $t \in T$. Il s'agit bien sûr d'un
cas particulier de la situation \ref{situation-numerical}, où
$I = \mathbf{Z} \to D(\mathcal{O}_X)$ est donné par
$n \mapsto \mathcal{L}^{\otimes n}$. Il résulte du lemme
\ref{lemma-finite-list-perfect-objects} qu'il s'agit d'un sous-champ ouvert
et fermé. Comme les fonctions
$n \mapsto
\chi(X_t, \mathcal{F}_t \otimes_{\mathcal{O}_{X_t}} \mathcal{L}_t^{\otimes n})$
sont toujours des polynômes numériques (Espaces sur les corps, lemme
\ref{spaces-over-fields-lemma-numerical-polynomial-from-euler}), on conclut
que
$$
\Cohstack_{X/B} = \coprod\nolimits_{P\text{ polynôme numérique}}
\Cohstack^P_{X/B}
$$
est une décomposition en union disjointe.
\end{example}







\section{Propriétés de Quot}
\label{section-quot}

\noindent
Soit $f : X \to B$ un morphisme d'espaces algébriques séparé et de
présentation finie. Soit $\mathcal{F}$ un $\mathcal{O}_X$-module
quasi-cohérent. Alors $\Quotfunctor_{\mathcal{F}/X/B}$ est un espace
algébrique. Si $\mathcal{F}$ est de présentation finie, alors
$\Quotfunctor_{\mathcal{F}/X/B} \to B$ est localement de présentation
finie. Voir Quotients, proposition \ref{quot-proposition-quot}.

\begin{lemma}
\label{lemma-quot-diagonal-closed}
La diagonale de $\Quotfunctor_{\mathcal{F}/X/B} \to B$ est une immersion
fermée. Si $\mathcal{F}$ est de type fini, la diagonale est une immersion
fermée de présentation finie.
\end{lemma}

\begin{proof}
Supposons donnés un schéma $T/B$ et deux quotients
$\mathcal{F}_T \to \mathcal{Q}_i$, $i = 1, 2$, correspondant à des points
à valeurs dans $T$ de $\Quotfunctor_{\mathcal{F}/X/B}$ au-dessus de $B$.
Notons $\mathcal{K}_1$ le noyau du premier et soit
$u : \mathcal{K}_1 \to \mathcal{Q}_2$ le morphisme composé.
D'après Platitude sur les espaces, lemme
\ref{spaces-flat-lemma-F-zero-closed-proper}, il existe un sous-espace fermé
de $T$ tel que $T' \to T$ se factorise par lui si et seulement si l'image
inverse $u_{T'}$ est nulle. Cela montre que la diagonale est une immersion
fermée. De plus, si $\mathcal{F}$ est de type fini, alors
$\mathcal{K}_1$ est de type fini (Modules sur les sites, lemme
\ref{sites-modules-lemma-kernel-surjection-finite-onto-finite-presentation})
et le même lemme montre que la diagonale est de présentation finie.
\end{proof}

\begin{lemma}
\label{lemma-quot-s-lfp}
Le morphisme $\Quotfunctor_{\mathcal{F}/X/B} \to B$ est séparé.
Si $\mathcal{F}$ est de présentation finie, il est aussi localement de
présentation finie.
\end{lemma}

\begin{proof}
Pour vérifier que $\Quotfunctor_{\mathcal{F}/X/B} \to B$ est séparé, il
faut montrer que sa diagonale est une immersion fermée. Cela résulte du
lemme \ref{lemma-quot-diagonal-closed}. La seconde assertion fait partie
de Quotients, proposition \ref{quot-proposition-quot}.
\end{proof}

\begin{lemma}
\label{lemma-quot-existence-part}
Supposons que $X \to B$ soit propre ainsi que de présentation finie et que
$\mathcal{F}$ soit quasi-cohérent de type fini.
Alors $\Quotfunctor_{\mathcal{F}/X/B} \to B$ satisfait la partie
d'existence du critère valuatif (Morphismes d'espaces, définition
\ref{spaces-morphisms-definition-valuative-criterion}).
\end{lemma}

\begin{proof}
Après changement de base, on se ramène immédiatement au problème suivant :
étant donnés un anneau de valuation $R$ de corps des fractions $K$, un
espace algébrique $X$ propre sur $R$, un $\mathcal{O}_X$-module
quasi-cohérent $\mathcal{F}$ de type fini et un quotient cohérent
$\mathcal{F}_K \to \mathcal{Q}_K$, montrer qu'il existe un quotient
$\mathcal{F} \to \mathcal{Q}$ où $\mathcal{Q}$ est un
$\mathcal{O}_X$-module de présentation finie, plat sur $R$, dont la fibre
générique est $\mathcal{Q}_K$.
Remarquons que, d'après Platitude sur les espaces, théorème
\ref{spaces-flat-theorem-finite-type-flat}, tout $\mathcal{O}_X$-module
quasi-cohérent $\mathcal{F}$ de type fini et plat sur $R$ est de
présentation finie. Nous établissons d'abord l'existence de $\mathcal{Q}$
localement sur les ouverts affines.

\medskip\noindent
Localement sur les ouverts affines, on aboutit au problème suivant :
soit $R \to A$ un homomorphisme d'anneaux de présentation finie, soit $M$
un $A$-module fini et soit $\varphi : M_K \to N_K$ un $A_K$-module
quotient. On peut alors considérer
$$
L = \{x \in M \mid \varphi(x \otimes 1) = 0 \}
$$
Le morphisme $M \to M/L$ définit un $A$-module quotient qui est sans
torsion comme $R$-module. Il est donc plat comme $R$-module (Compléments
d'algèbre, lemme
\ref{more-algebra-lemma-valuation-ring-torsion-free-flat}).
Puisque $M$ est fini comme $A$-module, il en va de même de $L$, et l'on
conclut que $L$ est de présentation finie comme $A$-module (par la
référence précédente). Manifestement, $M/L$ est l'unique quotient de ce
type tel que $(M/L)_K = N_K$.

\medskip\noindent
L'unicité dans la construction du paragraphe précédent garantit que ces
quotients se recollent et donnent le $\mathcal{Q}$ recherché. Donnons un
peu plus de détails. Choisissons un morphisme \'etale surjectif
$U \to X$, où $U$ est un schéma affine. La construction précédente donne
un quotient $\mathcal{F}|_U \to \mathcal{Q}_U$ qui est quasi-cohérent,
plat sur $R$ et redonne $\mathcal{Q}_K|U$ sur la fibre générique.
Comme $X$ est séparé, on voit que $U \times_X U$ est lui aussi un schéma
affine \'etale sur $X$.
Les quotients $\mathcal{F}|_{U \times_X U} \to
\text{pr}_1^*\mathcal{Q}_U$ et
$\mathcal{F}|_{U \times_X U} \to \text{pr}_2^*\mathcal{Q}_U$
coïncident alors par unicité de la construction. On peut donc descendre
$\mathcal{F}|_U \to \mathcal{Q}_U$ en une surjection
$\mathcal{F} \to \mathcal{Q}$, comme voulu (Propriétés des espaces,
proposition \ref{spaces-properties-proposition-quasi-coherent}).
\end{proof}

\begin{lemma}
\label{lemma-quot-functorial}
Soit $B$ un espace algébrique. Soit $\pi : X \to Y$ un morphisme
quasi-fini affine d'espaces algébriques séparés et de présentation finie sur
$B$. Soit $\mathcal{F}$ un $\mathcal{O}_X$-module quasi-cohérent.
Alors $\pi_*$ induit un morphisme
$\Quotfunctor_{\mathcal{F}/X/B} \to
\Quotfunctor_{\pi_*\mathcal{F}/Y/B}$.
\end{lemma}

\begin{proof}
Posons $\mathcal{G} = \pi_*\mathcal{F}$. Puisque $\pi$ est affine, pour
tout schéma $T$ sur $B$, on a
$\mathcal{G}_T = \pi_{T, *}\mathcal{F}_T$, d'après Cohomologie des
espaces, lemme \ref{spaces-cohomology-lemma-affine-base-change}.
De plus, $\pi_T$ est affine, donc $\pi_{T, *}$ est exact et transforme les
quotients en quotients. Remarquons qu'un quotient quasi-cohérent
$\mathcal{F}_T \to \mathcal{Q}$ définit un point de
$\Quotfunctor_{X/B}$ si et seulement si $\mathcal{Q}$ définit un objet de
$\Cohstack_{X/B}$ sur $T$ (et de même pour $\mathcal{G}$ et $Y$).
Comme nous avons vu au lemme \ref{lemma-coherent-functorial} que $\pi_*$
induit un morphisme $\Cohstack_{X/B} \to \Cohstack_{Y/B}$, on voit que si
$\mathcal{F}_T \to \mathcal{Q}$ appartient à
$\Quotfunctor_{\mathcal{F}/X/B}(T)$, alors
$\mathcal{G}_T \to \pi_{T, *}\mathcal{Q}$ appartient à
$\Quotfunctor_{\mathcal{G}/Y/B}(T)$.
\end{proof}

\begin{lemma}
\label{lemma-quot-open}
Soit $B$ un espace algébrique. Soit $\pi : X \to Y$ une immersion ouverte
affine d'espaces algébriques séparés et de présentation finie sur $B$.
Soit $\mathcal{F}$ un $\mathcal{O}_X$-module quasi-cohérent. Alors le
morphisme
$\Quotfunctor_{\mathcal{F}/X/B} \to
\Quotfunctor_{\pi_*\mathcal{F}/Y/B}$ du lemme
\ref{lemma-quot-functorial} est une immersion ouverte.
\end{lemma}

\begin{proof}
Omis. Indication : si $(\pi_*\mathcal{F})_T \to \mathcal{Q}$ est un
élément de $\Quotfunctor_{\pi_*\mathcal{F}/Y/B}(T)$ et si, pour
$t \in T$, on a $\text{Supp}(\mathcal{Q}_t) \subset |X_t|$, alors il en
va de même pour $t' \in T$ dans un voisinage de $t$.
\end{proof}

\begin{lemma}
\label{lemma-quot-better-open}
Soit $B$ un espace algébrique. Soit $j : X \to Y$ une immersion ouverte
d'espaces algébriques séparés et de présentation finie sur $B$.
Soit $\mathcal{G}$ un $\mathcal{O}_Y$-module quasi-cohérent et posons
$\mathcal{F} = j^*\mathcal{G}$. Il existe alors une immersion ouverte
$$
\Quotfunctor_{\mathcal{F}/X/B}
\longrightarrow
\Quotfunctor_{\mathcal{G}/Y/B}
$$
d'espaces algébriques sur $B$.
\end{lemma}

\begin{proof}
Si $\mathcal{F}_T \to \mathcal{Q}$ est un élément de
$\Quotfunctor_{\mathcal{F}/X/B}(T)$, on peut considérer
$\mathcal{G}_T \to j_{T, *}\mathcal{F}_T \to j_{T, *}\mathcal{Q}$.
L'examen des germes montre que ce morphisme est surjectif.
D'après le lemme \ref{lemma-coherent-functorial},
$j_{T, *}\mathcal{Q}$ est de présentation finie, plat sur $B$ et à support
propre sur $B$. On obtient ainsi un point à valeurs dans $T$ de
$\Quotfunctor_{\mathcal{G}/Y/B}$. Cela définit le morphisme du lemme.
Nous omettons la démonstration du fait qu'il s'agit d'une immersion ouverte.
Indication : si $\mathcal{G}_T \to \mathcal{Q}$ est un élément de
$\Quotfunctor_{\mathcal{G}/Y/B}(T)$ et si, pour $t \in T$, on a
$\text{Supp}(\mathcal{Q}_t) \subset |X_t|$, alors il en va de même pour
$t' \in T$ dans un voisinage de $t$.
\end{proof}

\begin{lemma}
\label{lemma-quot-closed}
Soit $B$ un espace algébrique. Soit $\pi : X \to Y$ une immersion fermée
d'espaces algébriques séparés et de présentation finie sur $B$.
Soit $\mathcal{F}$ un $\mathcal{O}_X$-module quasi-cohérent.
Alors le morphisme
$\Quotfunctor_{\mathcal{F}/X/B} \to
\Quotfunctor_{\pi_*\mathcal{F}/Y/B}$ du lemme
\ref{lemma-quot-functorial} est un isomorphisme.
\end{lemma}

\begin{proof}
Pour tout schéma $T$ sur $B$, le morphisme $\pi_T : X_T \to Y_T$ est une
immersion fermée. Le foncteur $\pi_{T, *}$ est donc une équivalence de
catégories entre $\QCoh(\mathcal{O}_{X_T})$ et la sous-catégorie pleine de
$\QCoh(\mathcal{O}_{Y_T})$ formée des modules quasi-cohérents annulés par
le faisceau d'idéaux de $X_T$ ; voir Morphismes d'espaces, lemme
\ref{spaces-morphisms-lemma-i-star-equivalence}.
Puisqu'un quotient de $(\pi_*\mathcal{F})_T$ est annulé par cet idéal, on
obtient la bijectivité de l'application
$\Quotfunctor_{\mathcal{F}/X/B}(T) \to
\Quotfunctor_{\pi_*\mathcal{F}/Y/B}(T)$
pour tout $T$, comme voulu.
\end{proof}

\begin{lemma}
\label{lemma-quot-quotient}
Soit $X \to B$ comme dans l'introduction de cette section. Soit
$\mathcal{F} \to \mathcal{G}$ une surjection de $\mathcal{O}_X$-modules
quasi-cohérents. Il existe alors une immersion fermée canonique
$\Quotfunctor_{\mathcal{G}/X/B} \to
\Quotfunctor_{\mathcal{F}/X/B}$.
\end{lemma}

\begin{proof}
Soit $\mathcal{K} = \Ker(\mathcal{F} \to \mathcal{G})$.
Par exactitude à droite des images inverses, on obtient une suite exacte
$\mathcal{K}_T \to \mathcal{F}_T \to \mathcal{G}_T \to 0$
pour tout schéma $T$ sur $B$.
En particulier, un quotient de $\mathcal{G}_T$ détermine un quotient de
$\mathcal{F}_T$, et l'on obtient notre transformation de foncteurs
$\Quotfunctor_{\mathcal{G}/X/B} \to
\Quotfunctor_{\mathcal{F}/X/B}$.
Cette transformation est une immersion fermée d'après Platitude sur les
espaces, lemme \ref{spaces-flat-lemma-F-zero-closed-proper}.
En effet, étant donné un élément $\mathcal{F}_T \to \mathcal{Q}$ de
$\Quotfunctor_{\mathcal{F}/X/B}(T)$, son image inverse à $T'/T$ appartient
à l'image de la transformation si et seulement si
$\mathcal{K}_{T'} \to \mathcal{Q}_{T'}$ est nul.
\end{proof}

\begin{remark}[Invariants numériques]
\label{remark-quot-numerical}
Soient $f : X \to B$ et $\mathcal{F}$ comme dans l'introduction de cette
section. Soit $I$ un ensemble et, pour $i \in I$, soit
$E_i \in D(\mathcal{O}_X)$ parfait. Soit $P : I \to \mathbf{Z}$ une
fonction. Rappelons que nous avons un morphisme
$$
\Quotfunctor_{\mathcal{F}/X/B} \longrightarrow \Cohstack_{X/B}
$$
qui envoie l'élément $\mathcal{F}_T \to \mathcal{Q}$ de
$\Quotfunctor_{\mathcal{F}/X/B}(T)$ sur l'objet $\mathcal{Q}$ de
$\Cohstack_{X/B}$ au-dessus de $T$ ; voir la démonstration de Quotients,
proposition \ref{quot-proposition-quot}. On peut donc former le diagramme
de produit fibré
$$
\xymatrix{
\Quotfunctor^P_{\mathcal{F}/X/B} \ar[r] \ar[d] &
\Cohstack^P_{X/B} \ar[d] \\
\Quotfunctor_{\mathcal{F}/X/B} \ar[r] &
\Cohstack_{X/B}
}
$$
Il s'agit du diagramme qui définit l'espace algébrique situé en haut à
gauche. La flèche verticale de gauche est une immersion fermée plate qui
est une immersion ouverte et fermée, par exemple si $I$ est fini, si $B$
est localement noethérien, ou si $I = \mathbf{Z}$ et
$E_i = \mathcal{L}^{\otimes i}$ pour un $\mathcal{O}_X$-module inversible
$\mathcal{L}$ (dans ce dernier cas, nous employons parfois la notation
$\Quotfunctor^{P, \mathcal{L}}_{\mathcal{F}/X/B}$).
Voir la situation \ref{situation-numerical}, les lemmes
\ref{lemma-open-P} et \ref{lemma-finite-list-perfect-objects}, ainsi que
l'exemple \ref{example-hilbert-polynomial}.
\end{remark}

\begin{lemma}
\label{lemma-quot-tensor-invertible}
Soient $f : X \to B$ et $\mathcal{F}$ comme dans l'introduction de cette
section. Soit $\mathcal{L}$ un $\mathcal{O}_X$-module inversible.
Le produit tensoriel par $\mathcal{L}$ définit alors un isomorphisme
$$
\Quotfunctor_{\mathcal{F}/X/B}
\longrightarrow
\Quotfunctor_{\mathcal{F} \otimes_{\mathcal{O}_X} \mathcal{L}/X/B}
$$
Étant donné un polynôme numérique $P(t)$, posons $P'(t) = P(t + 1)$ ;
cette application induit alors un isomorphisme
$\Quotfunctor^P_{\mathcal{F}/X/B}
\longrightarrow
\Quotfunctor^{P'}_{\mathcal{F} \otimes_{\mathcal{O}_X} \mathcal{L}/X/B}$
de sous-champs ouverts et fermés.
\end{lemma}

\begin{proof}
Posons $\mathcal{G} = \mathcal{F} \otimes_{\mathcal{O}_X} \mathcal{L}$.
Remarquons que
$\mathcal{G}_T = \mathcal{F}_T \otimes_{\mathcal{O}_{X_T}}
\mathcal{L}_T$.
Si $\mathcal{F}_T \to \mathcal{Q}$ est un élément de
$\Quotfunctor_{\mathcal{F}/X/B}(T)$, on l'envoie sur l'élément
$\mathcal{G}_T \to
\mathcal{Q} \otimes_{\mathcal{O}_{X_T}} \mathcal{L}_T$
de $\Quotfunctor_{\mathcal{F} \otimes_{\mathcal{O}_X}
\mathcal{L}/X/B}(T)$.
Cette construction est compatible aux images inverses et définit donc la
transformation de foncteurs voulue. Comme il existe une transformation
inverse évidente, c'est un isomorphisme. Nous omettons la démonstration de
la dernière assertion.
\end{proof}

\begin{lemma}
\label{lemma-quot-power-invertible}
Soient $f : X \to B$ et $\mathcal{F}$ comme dans l'introduction de cette
section. Soit $\mathcal{L}$ un $\mathcal{O}_X$-module inversible.
Alors
$$
\Quotfunctor^{P, \mathcal{L}}_{\mathcal{F}/X/B} =
\Quotfunctor^{P', \mathcal{L}^{\otimes n}}_{\mathcal{F}/X/B}
$$
où $P'(t) = P(nt)$.
\end{lemma}

\begin{proof}
Cela résulte immédiatement du développement des définitions.
\end{proof}

\section{Bornitude pour Quot}
\label{section-quot-bounded}

\noindent
Contrairement à ce qui se passe dans la situation classique, nous savons
déjà que le foncteur Quot est un espace algébrique, mais nous ignorons s'il
est jamais représenté par un espace algébrique de type fini.

\begin{lemma}
\label{lemma-quot-Pn}
Soient $n \geq 0$, $r \geq 1$ et $P \in \mathbf{Q}[t]$.
L'espace algébrique
$$
X = \Quotfunctor^P_{\mathcal{O}^{\oplus r}_{\mathbf{P}^n_\mathbf{Z}}/
\mathbf{P}^n_\mathbf{Z}/\mathbf{Z}}
$$
qui paramètre les quotients de
$\mathcal{O}_{\mathbf{P}^n_\mathbf{Z}}^{\oplus r}$ de polynôme de Hilbert
$P$ est propre sur $\Spec(\mathbf{Z})$.
\end{lemma}

\begin{proof}
Nous savons déjà que $X \to \Spec(\mathbf{Z})$ est séparé et localement de
présentation finie (lemme \ref{lemma-quot-s-lfp}).
Nous savons aussi que $X \to \Spec(\mathbf{Z})$ satisfait la partie
d'existence du critère valuatif ; voir le lemme
\ref{lemma-quot-existence-part}.
D'après le critère valuatif de propreté, il suffit de montrer que notre
espace Quot est quasi-compact ; voir Morphismes d'espaces, lemme
\ref{spaces-morphisms-lemma-characterize-proper}.
Il suffit donc de trouver un schéma quasi-compact $T$ et un morphisme
surjectif $T \to X$. Soit $m$ l'entier fourni par Variétés, lemme
\ref{varieties-lemma-bound-quotients-free}. Posons
$$
N = r{m + n \choose n} - P(m)
$$
Nous écrirons $\mathbf{P}^n$ pour
$\mathbf{P}^n_\mathbf{Z} = \text{Proj}(\mathbf{Z}[T_0, \ldots, T_n])$,
et les produits dépourvus d'indice seront pris sur $\Spec(\mathbf{Z})$.
L'idée de la démonstration consiste à construire une application
« universelle »
$$
\Psi :
\mathcal{O}_{T \times \mathbf{P}^n}(-m)^{\oplus N}
\longrightarrow
\mathcal{O}_{T \times \mathbf{P}^n}^{\oplus r}
$$
sur un schéma affine $T$, puis à montrer que tout point de $X$ correspond
au conoyau de cette application en un certain point de $T$.

\medskip\noindent
Définition de $T$ et de $\Psi$. Prenons $T = \Spec(A)$, où
$$
A = \mathbf{Z}[a_{i, j, E}]
$$
avec $i \in \{1, \ldots, r\}$, $j \in \{1, \ldots, N\}$ et où
$E = (e_0, \ldots, e_n)$ parcourt les multi-indices de degré total
$|E| = \sum_{k = 0, \ldots n} e_k = m$.
Nous définissons alors $\Psi$ comme l'application dont l'entrée de matrice
$(i, j)$ est l'application
$$
\sum\nolimits_{E = (e_0, \ldots, e_n)}
a_{i, j, E} T_0^{e_0} \ldots T_n^{e_n} :
\mathcal{O}_{T \times \mathbf{P}^n}(-m)
\longrightarrow
\mathcal{O}_{T \times \mathbf{P}^n}
$$
où la somme porte sur les $E$ ci-dessus ($i$ et $j$ étant bien entendu
fixés).

\medskip\noindent
Considérons le quotient $\mathcal{Q} = \Coker(\Psi)$ sur
$T \times \mathbf{P}^n$. D'après Compléments sur les morphismes, lemme
\ref{more-morphisms-lemma-generic-flatness-stratification}, il existe un
$t \geq 0$ et des sous-schémas fermés
$$
T = T_0 \supset T_1 \supset \ldots \supset T_t = \emptyset
$$
tels que l'image inverse $\mathcal{Q}_p$ de $\mathcal{Q}$ sur
$(T_p \setminus T_{p + 1}) \times \mathbf{P}^n$ soit plate sur
$T_p \setminus T_{p + 1}$. Remarquons que l'on a une suite exacte
$$
\mathcal{O}_{(T_p \setminus T_{p + 1}) \times \mathbf{P}^n}(-m)^{\oplus N}
\to
\mathcal{O}_{(T_p \setminus T_{p + 1}) \times \mathbf{P}^n}^{\oplus r}
\to
\mathcal{Q}_p
\to
0
$$
en prenant l'image inverse de la suite exacte qui définit
$\mathcal{Q} = \Coker(\Psi)$. Nous obtenons donc un morphisme
$$
\coprod (T_p \setminus T_{p + 1})
\longrightarrow
\Quotfunctor_{\mathcal{O}^{\oplus r}/\mathbf{P}/\mathbf{Z}}
\supset
\Quotfunctor^P_{\mathcal{O}^{\oplus r}/\mathbf{P}/\mathbf{Z}} = X
$$
Puisque le membre de gauche est un schéma noethérien et que l'inclusion du
membre de droite est ouverte, il suffit de montrer que tout point de $X$
appartient à l'image de ce morphisme.

\medskip\noindent
Soit $k$ un corps et soit $x \in X(k)$. Alors $x$ correspond à une
surjection $\mathcal{O}_{\mathbf{P}^n_k}^{\oplus r} \to \mathcal{F}$ de
$\mathcal{O}_{\mathbf{P}^n_k}$-modules cohérents telle que le polynôme de
Hilbert de $\mathcal{F}$ soit $P$. Considérons la suite exacte courte
$$
0 \to \mathcal{K} \to
\mathcal{O}_{\mathbf{P}^n_k}^{\oplus r} \to
\mathcal{F} \to 0
$$
D'après Variétés, lemme \ref{varieties-lemma-bound-quotients-free}, et le
choix de $m$, le faisceau $\mathcal{K}$ est $m$-régulier.
D'après Variétés, lemme
\ref{varieties-lemma-m-regular-globally-generated},
$\mathcal{K}(m)$ est engendré par ses sections globales.
D'après Variétés, lemme \ref{varieties-lemma-m-regular-up}, et la
définition de la $m$-régularité, on a
$H^i(\mathbf{P}^n_k, \mathcal{K}(m)) = 0$ pour $i > 0$.
Il vient donc
$$
\dim_k H^0(\mathbf{P}^n_k, \mathcal{K}(m)) =
\chi(\mathcal{K}(m)) =
\chi(\mathcal{O}_{\mathbf{P}^n_k}(m)^{\oplus r}) -
\chi(\mathcal{F}(m)) = N
$$
par le choix de $N$. Cela donne une surjection
$$
\mathcal{O}_{\mathbf{P}^n_k}^{\oplus N}
\longrightarrow
\mathcal{K}(m)
$$
Après torsion en sens inverse et en utilisant la suite exacte courte
ci-dessus, on voit que $\mathcal{F}$ est le conoyau d'une application
$$
\Psi_x :
\mathcal{O}_{\mathbf{P}^n_k}(-m)^{\oplus N}
\to
\mathcal{O}_{\mathbf{P}^n_k}^{\oplus r}
$$
Il existe un unique homomorphisme d'anneaux $\tau : A \to k$ tel que le
changement de base de $\Psi$ par le morphisme correspondant
$t = \Spec(\tau) : \Spec(k) \to T$ soit $\Psi_x$.
En effet, les entrées de la matrice $N \times r$ qui définit $\Psi_x$
sont des polynômes homogènes
$\sum \lambda_{i, j, E} T_0^{e_0} \ldots T_n^{e_n}$
de degré $m$ en $T_0, \ldots, T_n$, à coefficients
$\lambda_{i, j, E} \in k$, et l'on peut poser
$\tau(a_{i, j, E}) = \lambda_{i, j, E}$.
Alors $t \in T_p \setminus T_{p + 1}$ pour un certain $p$, et l'image de
$t$ par le morphisme ci-dessus est $x$, comme voulu.
\end{proof}

\begin{lemma}
\label{lemma-quot-Pn-over-base}
Soit $B$ un espace algébrique. Soit
$X = B \times \mathbf{P}^n_\mathbf{Z}$. Soit $\mathcal{L}$ l'image
inverse de $\mathcal{O}_{\mathbf{P}^n}(1)$ sur $X$.
Soit $\mathcal{F}$ un $\mathcal{O}_X$-module de présentation finie.
L'espace algébrique $\Quotfunctor^P_{\mathcal{F}/X/B}$ qui paramètre les
quotients de $\mathcal{F}$ de polynôme de Hilbert $P$ relativement à
$\mathcal{L}$ est propre sur $B$.
\end{lemma}

\begin{proof}
La question est locale pour la topologie \'etale sur $B$ ; voir Morphismes
d'espaces, lemme \ref{spaces-morphisms-lemma-proper-local}.
Nous pouvons donc supposer que $B$ est un schéma affine.
Dans ce cas, $\mathcal{L}$ est un module inversible ample sur $X$
(d'après Constructions, lemme \ref{constructions-lemma-ample-on-proj}, et la
définition des modules inversibles amples dans Propriétés, définition
\ref{properties-definition-ample}).
On peut donc trouver $r' \geq 0$ et $r \geq 0$, ainsi qu'une surjection
$$
\mathcal{O}_X^{\oplus r} \longrightarrow
\mathcal{F} \otimes_{\mathcal{O}_X} \mathcal{L}^{\otimes r'}
$$
d'après Propriétés, proposition
\ref{properties-proposition-characterize-ample}.
En vertu du lemme \ref{lemma-quot-tensor-invertible}, on peut remplacer
$\mathcal{F}$ par
$\mathcal{F} \otimes_{\mathcal{O}_X} \mathcal{L}^{\otimes r'}$
et $P(t)$ par $P(t + r')$.
Le lemme \ref{lemma-quot-quotient} fournit une immersion fermée
$$
\Quotfunctor^P_{\mathcal{F}/X/B}
\longrightarrow
\Quotfunctor^P_{\mathcal{O}_X^{\oplus r}/X/B}
$$
Comme nous avons montré que
$\Quotfunctor^P_{\mathcal{O}_X^{\oplus r}/X/B} \to B$ est propre au lemme
\ref{lemma-quot-Pn}, on conclut.
\end{proof}

\begin{lemma}
\label{lemma-quot-proper-over-base}
Soit $f : X \to B$ un morphisme propre et de présentation finie d'espaces
algébriques. Soit $\mathcal{F}$ un $\mathcal{O}_X$-module de présentation
finie. Soit $\mathcal{L}$ un $\mathcal{O}_X$-module inversible ample sur
$X/B$ ; voir Diviseurs sur les espaces, définition
\ref{spaces-divisors-definition-relatively-ample}.
L'espace algébrique $\Quotfunctor^P_{\mathcal{F}/X/B}$ qui paramètre les
quotients de $\mathcal{F}$ de polynôme de Hilbert $P$ relativement à
$\mathcal{L}$ est propre sur $B$.
\end{lemma}

\begin{proof}
La question est locale pour la topologie \'etale sur $B$ ; voir Morphismes
d'espaces, lemme \ref{spaces-morphisms-lemma-proper-local}.
Nous pouvons donc supposer que $B$ est un schéma affine.
Il existe alors une immersion fermée $i : X \to \mathbf{P}^n_B$ telle que
$i^*\mathcal{O}_{\mathbf{P}^n_B}(1) \cong \mathcal{L}^{\otimes d}$
pour un certain $d \geq 1$. Voir Morphismes, lemme
\ref{morphisms-lemma-quasi-projective-finite-type-over-S}.
Remplacer $\mathcal{L}$ par $\mathcal{L}^{\otimes d}$ et le polynôme
numérique $P(t)$ par $P(dt)$ ne modifie pas
$\Quotfunctor^P_{\mathcal{F}/X/B}$ ; certains détails sont omis.
Nous pouvons donc supposer
$\mathcal{L} = i^*\mathcal{O}_{\mathbf{P}^n_B}(1)$.
L'isomorphisme
$\Quotfunctor_{\mathcal{F}/X/B} \to
\Quotfunctor_{i_*\mathcal{F}/\mathbf{P}^n_B/B}$ du lemme
\ref{lemma-quot-closed} induit alors un isomorphisme
$\Quotfunctor^P_{\mathcal{F}/X/B} \cong
\Quotfunctor^P_{i_*\mathcal{F}/\mathbf{P}^n_B/B}$.
Comme $\Quotfunctor^P_{i_*\mathcal{F}/\mathbf{P}^n_B/B}$ est propre sur
$B$ d'après le lemme \ref{lemma-quot-Pn-over-base}, on conclut.
\end{proof}

\begin{lemma}
\label{lemma-quot-qc-over-base}
Soit $f : X \to B$ un morphisme séparé et de présentation finie d'espaces
algébriques. Soit $\mathcal{F}$ un $\mathcal{O}_X$-module de présentation
finie. Soit $\mathcal{L}$ un $\mathcal{O}_X$-module inversible ample sur
$X/B$ ; voir Diviseurs sur les espaces, définition
\ref{spaces-divisors-definition-relatively-ample}.
L'espace algébrique $\Quotfunctor^P_{\mathcal{F}/X/B}$ qui paramètre les
quotients de $\mathcal{F}$ de polynôme de Hilbert $P$ relativement à
$\mathcal{L}$ est séparé et de présentation finie sur $B$.
\end{lemma}

\begin{proof}
Nous avons déjà vu que $\Quotfunctor_{\mathcal{F}/X/B} \to B$ est séparé
et localement de présentation finie ; voir le lemme
\ref{lemma-quot-s-lfp}. Il suffit donc de montrer que le sous-espace ouvert
$\Quotfunctor^P_{\mathcal{F}/X/B}$ de la remarque
\ref{remark-quot-numerical} est quasi-compact sur $B$.

\medskip\noindent
La question est locale pour la topologie \'etale sur $B$
(Morphismes d'espaces, lemme
\ref{spaces-morphisms-lemma-quasi-compact-local}).
Nous pouvons donc supposer $B$ affine.

\medskip\noindent
Supposons $B = \Spec(\Lambda)$. Écrivons
$\Lambda = \colim \Lambda_i$ comme limite inductive de ses sous-algèbres
de type fini sur $\mathbf{Z}$. On peut alors trouver un $i$ et un système
$X_i, \mathcal{F}_i, \mathcal{L}_i$ comme dans le lemme sur
$B_i = \Spec(\Lambda_i)$, dont le changement de base à $B$ donne
$X, \mathcal{F}, \mathcal{L}$. Cela résulte de Limites d'espaces, lemmes
\ref{spaces-limits-lemma-descend-finite-presentation} (pour trouver $X_i$),
\ref{spaces-limits-lemma-descend-modules-finite-presentation} (pour trouver
$\mathcal{F}_i$), \ref{spaces-limits-lemma-descend-invertible-modules}
(pour trouver $\mathcal{L}_i$) et
\ref{spaces-limits-lemma-descend-separated} (pour rendre $X_i$ séparé).
Comme
$$
\Quotfunctor_{\mathcal{F}/X/B} = B \times_{B_i}
\Quotfunctor_{\mathcal{F}_i/X_i/B_i}
$$
et qu'il en va de même de $\Quotfunctor^P_{\mathcal{F}/X/B}$, on se
ramène au cas étudié dans le paragraphe suivant.

\medskip\noindent
Supposons $B$ affine et noethérien. Nous pouvons remplacer $\mathcal{L}$
par une puissance positive ; voir le lemme
\ref{lemma-quot-power-invertible}.
Nous pouvons ainsi supposer qu'il existe une immersion
$i : X \to \mathbf{P}^n_B$ telle que
$i^*\mathcal{O}_{\mathbf{P}^n}(1) = \mathcal{L}$.
D'après Morphismes, lemme
\ref{morphisms-lemma-quasi-compact-immersion}, il existe un sous-schéma
fermé $X' \subset \mathbf{P}^n_B$ tel que $i$ se factorise par une
immersion ouverte $j : X \to X'$.
D'après Propriétés, lemme
\ref{properties-lemma-lift-finite-presentation}, il existe un
$\mathcal{O}_{X'}$-module $\mathcal{G}$ de présentation finie tel que
$j^*\mathcal{G} = \mathcal{F}$. On obtient donc une immersion ouverte
$$
\Quotfunctor_{\mathcal{F}/X/B}
\longrightarrow
\Quotfunctor_{\mathcal{G}/X'/B}
$$
par le lemme \ref{lemma-quot-better-open}. Manifestement, cette immersion
ouverte envoie $\Quotfunctor^P_{\mathcal{F}/X/B}$ dans
$\Quotfunctor^P_{\mathcal{G}/X'/B}$. Or
$\Quotfunctor^P_{\mathcal{G}/X'/B}$ est propre sur $B$ d'après le lemme
\ref{lemma-quot-proper-over-base}.
Il est donc noethérien et, puisque tout ouvert d'un espace algébrique
noethérien est quasi-compact, la démonstration est achevée.
\end{proof}




\section{Propriétés du foncteur de Hilbert}
\label{section-hilb}

\noindent
Soit $f : X \to B$ un morphisme d'espaces algébriques séparé et de
présentation finie. Alors $\Hilbfunctor_{X/B}$ est un espace algébrique
localement de présentation finie sur $B$. Voir Quotients, proposition
\ref{quot-proposition-hilb}.

\begin{lemma}
\label{lemma-hilb-diagonal-closed}
La diagonale de $\Hilbfunctor_{X/B} \to B$ est une immersion fermée de
présentation finie.
\end{lemma}

\begin{proof}
Dans Quotients, lemme \ref{quot-lemma-hilb-is-quot}, nous avons vu que
$\Hilbfunctor_{X/B} = \Quotfunctor_{\mathcal{O}_X/X/B}$.
L'assertion résulte donc du lemme \ref{lemma-quot-diagonal-closed}.
\end{proof}

\begin{lemma}
\label{lemma-hilb-s-lfp}
Le morphisme $\Hilbfunctor_{X/B} \to B$ est séparé et localement de
présentation finie.
\end{lemma}

\begin{proof}
Pour vérifier que $\Hilbfunctor_{X/B} \to B$ est séparé, il faut montrer
que sa diagonale est une immersion fermée. Cela résulte du lemme
\ref{lemma-hilb-diagonal-closed}. La seconde assertion fait partie de
Quotients, proposition \ref{quot-proposition-hilb}.
\end{proof}

\begin{lemma}
\label{lemma-hilb-existence-part}
Supposons que $X \to B$ soit propre ainsi que de présentation finie.
Alors $\Hilbfunctor_{X/B} \to B$ satisfait la partie d'existence du critère
valuatif (Morphismes d'espaces, définition
\ref{spaces-morphisms-definition-valuative-criterion}).
\end{lemma}

\begin{proof}
Dans Quotients, lemme \ref{quot-lemma-hilb-is-quot}, nous avons vu que
$\Hilbfunctor_{X/B} = \Quotfunctor_{\mathcal{O}_X/X/B}$.
L'assertion résulte donc du lemme \ref{lemma-quot-existence-part}.
\end{proof}

\begin{lemma}
\label{lemma-hilb-open}
Soit $B$ un espace algébrique. Soit $\pi : X \to Y$ une immersion ouverte
d'espaces algébriques séparés et de présentation finie sur $B$.
Alors $\pi$ induit une immersion ouverte
$\Hilbfunctor_{X/B} \to \Hilbfunctor_{Y/B}$.
\end{lemma}

\begin{proof}
Omis. Indication : si $Z \subset X_T$ est un sous-schéma fermé propre sur
$T$, alors $Z$ est aussi fermé dans $Y_T$. On obtient ainsi la
transformation $\Hilbfunctor_{X/B} \to \Hilbfunctor_{Y/B}$.
Si $Z \subset Y_T$ est un élément de $\Hilbfunctor_{Y/B}(T)$ et si, pour
$t \in T$, on a $|Z_t| \subset |X_t|$, alors il en va de même pour
$t' \in T$ dans un voisinage de $t$.
\end{proof}

\begin{lemma}
\label{lemma-hilb-closed}
Soit $B$ un espace algébrique. Soit $\pi : X \to Y$ une immersion fermée
d'espaces algébriques séparés et de présentation finie sur $B$.
Alors $\pi$ induit une immersion fermée
$\Hilbfunctor_{X/B} \to \Hilbfunctor_{Y/B}$.
\end{lemma}

\begin{proof}
Puisque $\pi$ est une immersion fermée, il est immédiat qu'étant donné un
sous-schéma fermé $Z \subset X_T$, on peut voir $Z$ comme un sous-schéma
fermé de $X_T$. On obtient ainsi la transformation
$\Hilbfunctor_{X/B} \to \Hilbfunctor_{Y/B}$.
On voit immédiatement que cette transformation est un monomorphisme.
Pour prouver qu'elle est une immersion fermée, on peut appliquer le lemme
\ref{lemma-quot-quotient} à l'application
$\mathcal{O}_Y \to \mathcal{O}_X$ et aux identifications
$\Hilbfunctor_{X/B} = \Quotfunctor_{\mathcal{O}_X/X/B}$,
$\Hilbfunctor_{Y/B} = \Quotfunctor_{\mathcal{O}_Y/Y/B}$
de Quotients, lemme \ref{quot-lemma-hilb-is-quot}.
\end{proof}

\begin{remark}[Invariants numériques]
\label{remark-hilb-numerical}
Soit $f : X \to B$ comme dans l'introduction de cette section.
Soit $I$ un ensemble et, pour $i \in I$, soit
$E_i \in D(\mathcal{O}_X)$ parfait. Soit
$P : I \to \mathbf{Z}$ une fonction. Rappelons que
$\Hilbfunctor_{X/B} = \Quotfunctor_{\mathcal{O}_X/X/B}$ ; voir Quotients,
lemme \ref{quot-lemma-hilb-is-quot}. On peut donc définir
$$
\Hilbfunctor^P_{X/B} = \Quotfunctor^P_{\mathcal{O}_X/X/B}
$$
où $\Quotfunctor^P_{\mathcal{O}_X/X/B}$ est comme dans la remarque
\ref{remark-quot-numerical}. Le morphisme
$$
\Hilbfunctor^P_{X/B} \longrightarrow \Hilbfunctor_{X/B}
$$
est une immersion fermée plate qui est une immersion ouverte et fermée,
par exemple si $I$ est fini, si $B$ est localement noethérien, ou si
$I = \mathbf{Z}$ et $E_i = \mathcal{L}^{\otimes i}$ pour un
$\mathcal{O}_X$-module inversible $\mathcal{L}$. Dans ce dernier cas, nous
employons parfois la notation $\Hilbfunctor^{P, \mathcal{L}}_{X/B}$.
\end{remark}

\begin{lemma}
\label{lemma-hilb-proper-over-base}
Soit $f : X \to B$ un morphisme propre et de présentation finie d'espaces
algébriques. Soit $\mathcal{L}$ un $\mathcal{O}_X$-module inversible ample
sur $X/B$ ; voir Diviseurs sur les espaces, définition
\ref{spaces-divisors-definition-relatively-ample}.
L'espace algébrique $\Hilbfunctor^P_{X/B}$ qui paramètre les sous-schémas
fermés de polynôme de Hilbert $P$ relativement à $\mathcal{L}$ est propre
sur $B$.
\end{lemma}

\begin{proof}
Rappelons que
$\Hilbfunctor_{X/B} = \Quotfunctor_{\mathcal{O}_X/X/B}$ ; voir Quotients,
lemme \ref{quot-lemma-hilb-is-quot}.
Le présent lemme est donc une conséquence immédiate du lemme
\ref{lemma-quot-proper-over-base}.
\end{proof}

\begin{lemma}
\label{lemma-hilb-qc-over-base}
Soit $f : X \to B$ un morphisme séparé et de présentation finie d'espaces
algébriques. Soit $\mathcal{L}$ un $\mathcal{O}_X$-module inversible ample
sur $X/B$ ; voir Diviseurs sur les espaces, définition
\ref{spaces-divisors-definition-relatively-ample}.
L'espace algébrique $\Hilbfunctor^P_{X/B}$ qui paramètre les sous-schémas
fermés de polynôme de Hilbert $P$ relativement à $\mathcal{L}$ est séparé
et de présentation finie sur $B$.
\end{lemma}

\begin{proof}
Rappelons que
$\Hilbfunctor_{X/B} = \Quotfunctor_{\mathcal{O}_X/X/B}$ ; voir Quotients,
lemme \ref{quot-lemma-hilb-is-quot}.
Le présent lemme est donc une conséquence immédiate du lemme
\ref{lemma-quot-qc-over-base}.
\end{proof}








\section{Propriétés du champ de Picard}
\label{section-picard-stack}

\noindent
Soit $f : X \to B$ un morphisme d'espaces algébriques plat, propre et de
présentation finie. Alors le champ $\Picardstack_{X/B}$ qui paramètre les
faisceaux inversibles sur $X/B$ est algébrique ; voir Quotients,
proposition \ref{quot-proposition-pic}.

\begin{lemma}
\label{lemma-pic-diagonal-affine-fp}
La diagonale de $\Picardstack_{X/B}$ sur $B$ est affine et de présentation
finie.
\end{lemma}

\begin{proof}
Dans Quotients, lemme \ref{quot-lemma-picard-stack-open-in-coh}, nous avons
vu que $\Picardstack_{X/B}$ est un sous-champ ouvert de
$\Cohstack_{X/B}$. L'assertion résulte donc du lemme
\ref{lemma-coherent-diagonal-affine-fp}.
\end{proof}

\begin{lemma}
\label{lemma-pic-qs-lfp}
Le morphisme $\Picardstack_{X/B} \to B$ est quasi-séparé et localement de
présentation finie.
\end{lemma}

\begin{proof}
Dans Quotients, lemme \ref{quot-lemma-picard-stack-open-in-coh}, nous avons
vu que $\Picardstack_{X/B}$ est un sous-champ ouvert de
$\Cohstack_{X/B}$. L'assertion résulte donc du lemme
\ref{lemma-coherent-qs-lfp}.
\end{proof}

\begin{lemma}
\label{lemma-pic-existence-part}
Supposons en outre que $X \to B$ soit lisse, en plus d'être propre.
Alors $\Picardstack_{X/B} \to B$ satisfait la partie d'existence du critère
valuatif (Morphismes de champs, définition
\ref{stacks-morphisms-definition-existence}).
\end{lemma}

\begin{proof}
Après changement de base, on se ramène immédiatement au problème suivant :
étant donnés un anneau de valuation $R$ de corps des fractions $K$, un
espace algébrique $X$ propre et lisse sur $R$ et un
$\mathcal{O}_{X_K}$-module inversible $\mathcal{L}_K$, montrer qu'il existe
un $\mathcal{O}_X$-module inversible $\mathcal{L}$ dont la fibre générique
est $\mathcal{L}_K$.
Remarquons que $X_K$ est noethérien, séparé et régulier (utiliser Morphismes
d'espaces, lemme
\ref{spaces-morphisms-lemma-finite-presentation-noetherian}, et Espaces sur
les corps, lemme \ref{spaces-over-fields-lemma-smooth-regular}).
On peut donc écrire $\mathcal{L}_K$ comme la différence, dans le groupe de
Picard, de $\mathcal{O}_{X_K}(D_K)$ et
$\mathcal{O}_{X_K}(D'_K)$, pour deux diviseurs de Cartier effectifs
$D_K, D'_K$ dans $X_K$ ; voir Diviseurs sur les espaces, lemme
\ref{spaces-divisors-lemma-Noetherian-regular-separated-pic-effective-Cartier}.
Enfin, on sait que $D_K$ et $D'_K$ sont les restrictions de diviseurs de
Cartier effectifs $D, D' \subset X$ ; voir Diviseurs sur les espaces,
lemme
\ref{spaces-divisors-lemma-smooth-over-valuation-ring-effective-Cartier}.
\end{proof}

\begin{lemma}
\label{lemma-pic-inertia}
Supposons $f_{T, *}\mathcal{O}_{X_T} \cong \mathcal{O}_T$ pour tout schéma
$T$ sur $B$. Alors le champ d'inertie de $\Picardstack_{X/B}$ est égal à
$\mathbf{G}_m \times \Picardstack_{X/B}$.
\end{lemma}

\begin{proof}
Ceci est expliqué dans Exemples de champs, exemple
\ref{examples-stacks-example-inertia-stack-of-picard}.
\end{proof}

\begin{lemma}
\label{lemma-pic-curves-smooth}
Supposons en outre que $f : X \to B$ soit de dimension relative
$\leq 1$, en plus des autres hypothèses de cette section. Alors
$\Picardstack_{X/B} \to B$ est lisse.
\end{lemma}

\begin{proof}
Nous savons déjà que $\Picardstack_{X/B} \to B$ est localement de
présentation finie ; voir le lemme \ref{lemma-pic-qs-lfp}.
Il suffit donc de montrer que $\Picardstack_{X/B} \to B$ est formellement
lisse ; voir Compléments sur les morphismes de champs, lemme
\ref{stacks-more-morphisms-lemma-smooth-formally-smooth}.
Après changement de base, on se ramène immédiatement au problème suivant :
étant donné un épaississement du premier ordre $T \subset T'$ de schémas
affines, étant donné $X' \to T'$ propre, plat, de présentation finie et de
dimension relative $\leq 1$, et étant donné sur
$X = T \times_{T'} X'$ un $\mathcal{O}_X$-module inversible
$\mathcal{L}$, prouver qu'il existe un $\mathcal{O}_{X'}$-module inversible
$\mathcal{L}'$ dont la restriction à $X$ est $\mathcal{L}$.
Puisque $T \subset T'$ est un épaississement du premier ordre, il en va de
même de $X \subset X'$ ; voir Compléments sur les morphismes d'espaces,
lemme \ref{spaces-more-morphisms-lemma-base-change-thickening}.
D'après Compléments sur les morphismes d'espaces, lemme
\ref{spaces-more-morphisms-lemma-picard-group-first-order-thickening}, il
suffit de montrer que $H^2(X, \mathcal{I}) = 0$, où $\mathcal{I}$ est
l'idéal quasi-cohérent qui définit $X$ dans $X'$.
Notons $f : X \to T$ le morphisme structural.
D'après Cohomologie des espaces, lemme
\ref{spaces-cohomology-lemma-higher-direct-images-zero-above-dimension-fibre},
on a $R^pf_*\mathcal{I} = 0$ pour $p > 1$.
La nullité recherchée résulte donc de Cohomologie des espaces, lemme
\ref{spaces-cohomology-lemma-quasi-coherence-higher-direct-images-application}
(c'est ici que nous utilisons enfin le fait que $T$ est affine).
\end{proof}




\section{Propriétés du foncteur de Picard}
\label{section-picard-functor}

\noindent
Soit $f : X \to B$ un morphisme d'espaces algébriques plat, propre et de
présentation finie, tel que, de plus, pour tout $T/B$, l'application
canonique
$$
\mathcal{O}_T \longrightarrow f_{T, *}\mathcal{O}_{X_T}
$$
soit un isomorphisme. Alors le foncteur de Picard
$\Picardfunctor_{X/B}$ est un espace algébrique ; voir Quotients,
proposition \ref{quot-proposition-pic-functor}.
Il existe un lien étroit avec le champ de Picard.

\begin{lemma}
\label{lemma-pic-gerbe-over-pic-functor}
Le morphisme $\Picardstack_{X/B} \to \Picardfunctor_{X/B}$ fait du champ
de Picard une gerbe au-dessus du foncteur de Picard.
\end{lemma}

\begin{proof}
La définition du fait que
$\Picardstack_{X/B} \to \Picardfunctor_{X/B}$ soit une gerbe est donnée
dans Morphismes de champs, définition
\ref{stacks-morphisms-definition-gerbe}, qui renvoie à son tour à Champs,
définition \ref{stacks-definition-gerbe-over-stack-in-groupoids}.
Pour le prouver, nous allons vérifier les conditions (2)(a) et (2)(b) de
Champs, lemme \ref{stacks-lemma-when-gerbe}. Cela résulte immédiatement de
Quotients, lemme \ref{quot-lemma-pic-over-pic} ; voici une explication
détaillée.

\medskip\noindent
Condition (2)(a).
Supposons que $\xi \in \Picardfunctor_{X/B}(U)$ pour un schéma $U$ sur
$B$. Comme $\Picardfunctor_{X/B}$ est le faisceau associé pour la topologie
fppf à la règle $T \mapsto \Pic(X_T)$ sur les schémas au-dessus de $B$
(Quotients, situation \ref{quot-situation-pic}), il existe un recouvrement
fppf $\{U_i \to U\}$ tel que $\xi|_{U_i}$ corresponde à un module
inversible $\mathcal{L}_i$ sur $X_{U_i}$.
Alors $(U_i \to B, \mathcal{L}_i)$ est un objet de
$\Picardstack_{X/B}$ sur $U_i$ qui s'envoie sur $\xi|_{U_i}$.

\medskip\noindent
Condition (2)(b). Supposons que $U$ soit un schéma sur $B$ et que
$\mathcal{L}, \mathcal{N}$ soient des modules inversibles sur $X_U$ qui
s'envoient sur le même élément de $\Picardfunctor_{X/B}(U)$.
Il existe alors un recouvrement fppf $\{U_i \to U\}$ tel que
$\mathcal{L}|_{X_{U_i}}$ soit isomorphe à
$\mathcal{N}|_{X_{U_i}}$. On obtient ainsi les isomorphismes recherchés
entre
$(U \to B, \mathcal{L})|_{U_i} \to
(U \to B, \mathcal{N})|_{U_i}$.
\end{proof}

\begin{lemma}
\label{lemma-pic-functor-diagonal-qc-immersion}
La diagonale de $\Picardfunctor_{X/B}$ sur $B$ est une immersion
quasi-compacte.
\end{lemma}

\begin{proof}
La diagonale est une immersion d'après Quotients, lemme
\ref{quot-lemma-diagonal-pic}. Pour conclure, montrons qu'elle est
quasi-compacte. La diagonale de $\Picardstack_{X/B}$ est quasi-compacte
d'après le lemme \ref{lemma-pic-diagonal-affine-fp}, et
$\Picardstack_{X/B}$ est une gerbe au-dessus de
$\Picardfunctor_{X/B}$ d'après le lemme
\ref{lemma-pic-gerbe-over-pic-functor}.
On conclut par Morphismes de champs, lemme
\ref{stacks-morphisms-lemma-gerbe-diagonal-quasi-compact}.
\end{proof}

\begin{lemma}
\label{lemma-pic-functor-qs-lfp}
Le morphisme $\Picardfunctor_{X/B} \to B$ est quasi-séparé et localement
de présentation finie.
\end{lemma}

\begin{proof}
Pour vérifier que $\Picardfunctor_{X/B} \to B$ est quasi-séparé, il faut
montrer que sa diagonale est quasi-compacte. Cela résulte immédiatement du
lemme \ref{lemma-pic-functor-diagonal-qc-immersion}.
Comme le morphisme
$\Picardstack_{X/B} \to \Picardfunctor_{X/B}$ est surjectif, plat et
localement de présentation finie (d'après le lemme
\ref{lemma-pic-gerbe-over-pic-functor} et Morphismes de champs, lemme
\ref{stacks-morphisms-lemma-gerbe-fppf}), il suffit de prouver que
$\Picardstack_{X/B} \to B$ est localement de présentation finie ; voir
Morphismes de champs, lemme
\ref{stacks-morphisms-lemma-flat-finite-presentation-permanence}.
Cela résulte du lemme \ref{lemma-pic-qs-lfp}.
\end{proof}

\begin{lemma}
\label{lemma-pic-functor-uniqueness-part}
Supposons en outre que les fibres géométriques de $X \to B$ soient
intègres, en plus des autres hypothèses de cette section.
Alors $\Picardfunctor_{X/B} \to B$ est séparé.
\end{lemma}

\begin{proof}
Puisque $\Picardfunctor_{X/B} \to B$ est quasi-séparé, il suffit de
vérifier la partie d'unicité du critère valuatif ; voir Morphismes
d'espaces, lemme
\ref{spaces-morphisms-lemma-valuative-criterion-separatedness}.
On se ramène immédiatement au problème suivant : étant donnés
\begin{enumerate}
\item un anneau de valuation $R$ de corps des fractions $K$,
\item un espace algébrique $X$ propre et plat sur $R$, à fibre géométrique
intègre,
\item un élément $a \in \Picardfunctor_{X/R}(R)$ tel que
$a|_{\Spec(K)} = 0$,
\end{enumerate}
il faut prouver que $a = 0$. Appliquons Morphismes de champs, lemme
\ref{stacks-morphisms-lemma-lift-valuation-ring-through-flat-morphism},
au morphisme plat surjectif
$\Picardstack_{X/R} \to \Picardfunctor_{X/R}$
(surjectif et plat d'après le lemme
\ref{lemma-pic-gerbe-over-pic-functor} et Morphismes de champs, lemme
\ref{stacks-morphisms-lemma-gerbe-fppf}). Après avoir remplacé $R$ par une
extension, on peut supposer que $a$ est donné par un
$\mathcal{O}_X$-module inversible $\mathcal{L}$.
Comme $a|_{\Spec(K)} = 0$, on obtient
$\mathcal{L}_K \cong \mathcal{O}_{X_K}$ d'après Quotients, lemme
\ref{quot-lemma-flat-geometrically-connected-fibres}.

\medskip\noindent
Notons $f : X \to \Spec(R)$ le morphisme structural.
Soient $\eta, 0 \in \Spec(R)$ le point générique et le point fermé.
Considérons les complexes parfaits
$K = Rf_*\mathcal{L}$ et $M = Rf_*(\mathcal{L}^{\otimes -1})$
sur $\Spec(R)$ ; voir Catégories dérivées des espaces, lemme
\ref{spaces-perfect-lemma-flat-proper-perfect-direct-image-general}.
Considérons les fonctions
$\beta_{K, i}, \beta_{M, i} : \Spec(R) \to \mathbf{Z}$
de Catégories dérivées des espaces, lemme
\ref{spaces-perfect-lemma-jump-loci}, associées à $K$ et à $M$.
Comme la formation de $K$ et de $M$ commute au changement de base (voir le
lemme cité ci-dessus), on obtient
$\beta_{K, 0}(\eta) = \beta_{M, 0}(\eta) = 1$ d'après Espaces sur les
corps, lemme
\ref{spaces-over-fields-lemma-proper-geometrically-reduced-global-sections},
et notre hypothèse sur les fibres de $f$.
Comme ces fonctions sont semi-continues supérieurement, on obtient
$\beta_{K, 0}(0) \geq 1$ et $\beta_{M, 0}(0) \geq 1$.
D'après Espaces sur les corps, lemme
\ref{spaces-over-fields-lemma-characterize-trivial-pic-integral}, on
conclut que la restriction de $\mathcal{L}$ à la fibre spéciale $X_0$ est
triviale. Cela donne à son tour
$\beta_{K, 0}(0) = \beta_{M, 0} = 1$, comme ci-dessus.
D'après Compléments d'algèbre, lemme
\ref{more-algebra-lemma-lift-pseudo-coherent-from-residue-field}, on peut
alors représenter $K$ par un complexe de la forme
$$
\ldots \to 0 \to R \to R^{\oplus \beta_{K, 1}(0)} \to
R^{\oplus \beta_{K, 2}(0)} \to \ldots
$$
Or $R \to R^{\oplus \beta_{K, 1}(0)}$ est nul parce que
$\beta_{K, 0}(\eta) = 1$. Autrement dit,
$K = R \oplus \tau_{\geq 1}(K)$ dans $D(R)$, où
$\tau_{\geq 1}(K)$ est d'amplitude de Tor dans $[1, b]$ pour un certain
$b \in \mathbf{Z}$.
Il existe donc une section globale $s \in H^0(X, \mathcal{L})$ dont la
restriction $s_0$ à $X_0$ ne s'annule nulle part (de nouveau parce que la
formation de $K$ commute au changement de base).
Alors $s : \mathcal{O}_X \to \mathcal{L}$ est une application de faisceaux
inversibles dont la restriction à $X_0$ est un isomorphisme ; c'est donc
un isomorphisme, comme voulu.
\end{proof}

\begin{lemma}
\label{lemma-pic-functor-curves-smooth}
Supposons en outre que $f : X \to B$ soit de dimension relative
$\leq 1$, en plus des autres hypothèses de cette section. Alors
$\Picardfunctor_{X/B} \to B$ est lisse.
\end{lemma}

\begin{proof}
D'après le lemme \ref{lemma-pic-curves-smooth}, nous savons que
$\Picardstack_{X/B} \to B$ est lisse. Le morphisme
$\Picardstack_{X/B} \to \Picardfunctor_{X/B}$ est surjectif et lisse en
combinant le lemme \ref{lemma-pic-gerbe-over-pic-functor} et Morphismes de
champs, lemme \ref{stacks-morphisms-lemma-gerbe-smooth}.
Ainsi, si $U$ est un schéma et si
$U \to \Picardstack_{X/B}$ est surjectif et lisse, alors
$U \to \Picardfunctor_{X/B}$ est surjectif et lisse, et
$U \to B$ est surjectif et lisse (car ces propriétés sont stables par
composition). Le morphisme $\Picardfunctor_{X/B} \to B$ est donc lisse,
par exemple d'après Descente sur les espaces, lemme
\ref{spaces-descent-lemma-syntomic-smooth-etale-permanence}.
\end{proof}






\section{Propriétés des morphismes relatifs}
\label{section-relative-morphisms}

\noindent
Soit $B$ un espace algébrique. Soient $X$ et $Y$ des espaces algébriques
sur $B$ tels que $Y \to B$ soit plat, propre et de présentation finie, et
que $X \to B$ soit séparé et de présentation finie.
Alors le foncteur $\mathit{Mor}_B(Y, X)$ des morphismes relatifs est un
espace algébrique localement de présentation finie sur $B$.
Voir Quotients, proposition \ref{quot-proposition-Mor}.

\begin{lemma}
\label{lemma-Mor-diagonal-closed}
La diagonale de $\mathit{Mor}_B(Y, X) \to B$ est une immersion fermée de
présentation finie.
\end{lemma}

\begin{proof}
Il existe une immersion ouverte
$\mathit{Mor}_B(Y, X) \to \Hilbfunctor_{Y \times_B X/B}$ ; voir
Quotients, lemme \ref{quot-lemma-Mor-into-Hilb-open}.
Le lemme résulte donc du lemme \ref{lemma-hilb-diagonal-closed}.
\end{proof}

\begin{lemma}
\label{lemma-Mor-s-lfp}
Le morphisme $\mathit{Mor}_B(Y, X) \to B$ est séparé et localement de
présentation finie.
\end{lemma}

\begin{proof}
Pour vérifier que $\mathit{Mor}_B(Y, X) \to B$ est séparé, il faut montrer
que sa diagonale est une immersion fermée. Cela résulte du lemme
\ref{lemma-Mor-diagonal-closed}. La seconde assertion fait partie de
Quotients, proposition \ref{quot-proposition-Mor}.
\end{proof}

\begin{lemma}
\label{lemma-Isom-in-Mor}
Avec $B, X, Y$ comme dans l'introduction de cette section, supposons en
outre que $X \to B$ soit propre. Alors le sous-foncteur
$\mathit{Isom}_B(Y, X) \subset \mathit{Mor}_B(Y, X)$ des isomorphismes est
un sous-espace ouvert.
\end{lemma}

\begin{proof}
Cela résulte immédiatement de Compléments sur les morphismes d'espaces,
lemme \ref{spaces-more-morphisms-lemma-where-isomorphism}.
\end{proof}

\begin{remark}[Invariants numériques]
\label{remark-Mor-numerical}
Soient $B, X, Y$ comme dans l'introduction de cette section.
Soit $I$ un ensemble et, pour $i \in I$, soit
$E_i \in D(\mathcal{O}_{Y \times_B X})$ parfait.
Soit $P : I \to \mathbf{Z}$ une fonction. Rappelons que
$$
\mathit{Mor}_B(Y, X) \subset
\Hilbfunctor_{Y \times_B X/B}
$$
est un sous-espace ouvert ; voir Quotients, lemme
\ref{quot-lemma-Mor-into-Hilb-open}. On peut donc définir
$$
\mathit{Mor}^P_B(Y, X) =
\mathit{Mor}_B(Y, X) \cap \Hilbfunctor^P_{Y \times_B X/B}
$$
où $\Hilbfunctor^P_{Y \times_B X/B}$ est comme dans la remarque
\ref{remark-hilb-numerical}. Le morphisme
$$
\mathit{Mor}^P_B(Y, X) \longrightarrow \mathit{Mor}_B(Y, X)
$$
est une immersion fermée plate qui est une immersion ouverte et fermée,
par exemple si $I$ est fini, si $B$ est localement noethérien, ou si
$I = \mathbf{Z}$, $E_i = \mathcal{L}^{\otimes i}$ pour un
$\mathcal{O}_{Y \times_B X}$-module inversible $\mathcal{L}$.
Dans ce dernier cas, nous employons parfois la notation
$\mathit{Mor}^{P, \mathcal{L}}_B(Y, X)$.
\end{remark}

\begin{lemma}
\label{lemma-Mor-qc-over-base}
Avec $B, X, Y$ comme dans l'introduction de cette section, soit
$\mathcal{L}$ ample sur $X/B$ et soit $\mathcal{N}$ ample sur $Y/B$.
Voir Diviseurs sur les espaces, définition
\ref{spaces-divisors-definition-relatively-ample}.
Soit $P$ un polynôme numérique. Alors
$$
\mathit{Mor}^{P, \mathcal{M}}_B(Y, X) \longrightarrow B
$$
est séparé et de présentation finie, où
$\mathcal{M} = \text{pr}_1^*\mathcal{N}
\otimes_{\mathcal{O}_{Y \times_B X}} \text{pr}_2^*\mathcal{L}$.
\end{lemma}

\begin{proof}
D'après le lemme \ref{lemma-Mor-s-lfp}, le morphisme
$\mathit{Mor}_B(Y, X) \to B$ est séparé et localement de présentation
finie. Il suffit donc de montrer que le sous-espace ouvert et fermé
$\mathit{Mor}^{P, \mathcal{M}}_B(Y, X)$ de la remarque
\ref{remark-Mor-numerical} est quasi-compact sur $B$.

\medskip\noindent
La question est locale pour la topologie \'etale sur $B$
(Morphismes d'espaces, lemme
\ref{spaces-morphisms-lemma-quasi-compact-local}).
Nous pouvons donc supposer $B$ affine.

\medskip\noindent
Supposons $B = \Spec(\Lambda)$. Notons que $X$ et $Y$ sont des schémas et
que $\mathcal{L}$ et $\mathcal{N}$ sont des faisceaux inversibles amples
sur $X$ et $Y$ (cela résulte immédiatement des définitions).
Écrivons $\Lambda = \colim \Lambda_i$ comme limite inductive de ses
sous-algèbres de type fini sur $\mathbf{Z}$. On peut alors trouver un $i$
et un système $X_i, Y_i, \mathcal{L}_i, \mathcal{N}_i$ comme dans le lemme
sur $B_i = \Spec(\Lambda_i)$, dont le changement de base à $B$ donne
$X, Y, \mathcal{L}, \mathcal{N}$. Cela résulte de Limites, lemmes
\ref{limits-lemma-descend-finite-presentation} (pour trouver $X_i$, $Y_i$),
\ref{limits-lemma-descend-invertible-modules} (pour trouver
$\mathcal{L}_i$, $\mathcal{N}_i$),
\ref{limits-lemma-descend-separated-finite-presentation} (pour rendre
$X_i \to B_i$ séparé), \ref{limits-lemma-eventually-proper} (pour rendre
$Y_i \to B_i$ propre) et \ref{limits-lemma-limit-ample} (pour rendre
$\mathcal{L}_i$, $\mathcal{N}_i$ amples).
Comme
$$
\mathit{Mor}_B(Y, X) = B \times_{B_i} \mathit{Mor}_{B_i}(Y_i, X_i)
$$
et qu'il en va de même de $\mathit{Mor}^P_B(Y, X)$, on se ramène au cas
étudié dans le paragraphe suivant.

\medskip\noindent
Supposons $B$ affine et noethérien. D'après Propriétés, lemme
\ref{properties-lemma-ample-on-product}, $\mathcal{M}$ est ample.
D'après le lemme \ref{lemma-hilb-qc-over-base},
$\Hilbfunctor^{P, \mathcal{M}}_{Y \times_B X/B}$ est de présentation finie
sur $B$, donc noethérien. Par construction,
$$
\mathit{Mor}^{P, \mathcal{M}}_B(Y, X) =
\mathit{Mor}_B(Y, X) \cap
\Hilbfunctor^{P, \mathcal{M}}_{Y \times_B X/B}
$$
est un sous-espace ouvert de
$\Hilbfunctor^{P, \mathcal{M}}_{Y \times_B X/B}$, donc quasi-compact
(puisque tout ouvert d'un espace algébrique noethérien est quasi-compact).
\end{proof}






\section{Propriétés du champ des schémas propres polarisés}
\label{section-polarized}

\noindent
Dans cette section, nous étudions des propriétés du champ de modules
$$
\Polarizedstack \longrightarrow \Spec(\mathbf{Z})
$$
dont la catégorie des sections sur un schéma $S$ est la catégorie des
schémas propres, plats et de présentation finie sur $S$, munis d'un
faisceau inversible relativement ample. C'est un champ algébrique d'après
Quotients, théorème \ref{quot-theorem-polarized-algebraic}.

\begin{lemma}
\label{lemma-polarized-diagonal-separated-fp}
La diagonale de $\Polarizedstack$ est séparée et de présentation finie.
\end{lemma}

\begin{proof}
Rappelons que $\Polarizedstack$ est un champ algébrique qui préserve les
limites ; voir Quotients, lemme \ref{quot-lemma-polarized-limits}.
D'après Limites de champs, lemme
\ref{stacks-limits-lemma-limit-preserving-diagonal}, cela implique que
$\Delta : \Polarizedstack \to \Polarizedstack \times \Polarizedstack$
préserve les limites. Ainsi, $\Delta$ est localement de présentation finie
d'après Limites de champs, proposition
\ref{stacks-limits-proposition-characterize-locally-finite-presentation}.

\medskip\noindent
Montrons que $\Delta$ est séparée. Pour cela, il suffit de montrer qu'étant
donnés un schéma affine $U$ et deux objets
$\upsilon = (Y, \mathcal{N})$ et $\chi = (X, \mathcal{L})$
de $\Polarizedstack$ sur $U$, l'espace algébrique
$$
\mathit{Isom}_{\Polarizedstack}(\upsilon, \chi)
$$
est séparé. La règle qui, à un isomorphisme
$\upsilon_T \to \chi_T$, associe l'isomorphisme sous-jacent
$Y_T \to X_T$ définit un morphisme
$$
\mathit{Isom}_{\Polarizedstack}(\upsilon, \chi)
\longrightarrow
\mathit{Isom}_U(Y, X)
$$
Puisque nous avons vu aux lemmes \ref{lemma-Mor-s-lfp} et
\ref{lemma-Isom-in-Mor} que le but est un espace algébrique séparé, il
suffit de prouver que ce morphisme est séparé.
Étant donné un isomorphisme $f : Y_T \to X_T$ sur un schéma $T/U$, on a
manifestement
$$
\mathit{Isom}_{\Polarizedstack}(\upsilon, \chi)
\times_{\mathit{Isom}_U(Y, X), [f]} T
=
\mathit{Isom}(\mathcal{N}_T, f^*\mathcal{L}_T)
$$
Ici, $[f] : T \to \mathit{Isom}_U(Y, X)$ désigne le point à valeurs dans
$T$ correspondant à $f$, et
$\mathit{Isom}(\mathcal{N}_T, f^*\mathcal{L}_T)$ est l'espace algébrique
étudié à la section \ref{section-hom-isom}.
Comme cet espace algébrique est affine sur $U$, l'assertion implique que
$\Delta$ est séparée.

\medskip\noindent
Pour achever la démonstration, montrons que $\Delta$ est quasi-compacte.
Comme $\Delta$ est représentable par des espaces algébriques, il suffit de
vérifier que le changement de base de $\Delta$ par un morphisme lisse
surjectif $U \to \Polarizedstack \times \Polarizedstack$ est quasi-compact
(voir par exemple Propriétés des champs, lemme
\ref{stacks-properties-lemma-check-property-covering}).
Nous pouvons supposer que $U = \coprod U_i$ est une union disjointe
d'ouverts affines. Comme $\Polarizedstack$ préserve les limites (voir
ci-dessus), on voit que
$\Polarizedstack \to \Spec(\mathbf{Z})$ est localement de présentation
finie, donc que $U_i \to \Spec(\mathbf{Z})$ est localement de présentation
finie (Limites de champs, proposition
\ref{stacks-limits-proposition-characterize-locally-finite-presentation},
et Morphismes de champs, lemmes
\ref{stacks-morphisms-lemma-composition-finite-presentation} et
\ref{stacks-morphisms-lemma-smooth-locally-finite-presentation}).
En particulier, $U_i$ est affine noethérien. On se ramène ainsi au cas
étudié dans le paragraphe suivant.

\medskip\noindent
Dans ce paragraphe, étant donnés un schéma affine noethérien $U$ et deux
objets $\upsilon = (Y, \mathcal{N})$ et
$\chi = (X, \mathcal{L})$ de $\Polarizedstack$ sur $U$, montrons que
l'espace algébrique
$$
\mathit{Isom}_{\Polarizedstack}(\upsilon, \chi)
$$
est quasi-compact. Comme les composantes connexes de $U$ sont ouvertes et
fermées, on peut remplacer $U$ par celles-ci. Nous pouvons donc supposer,
et supposerons, que $U$ est connexe.
Soit $u \in U$ un point. Soit $P$ le polynôme de Hilbert
$n \mapsto \chi(Y_u, \mathcal{N}_u^{\otimes n})$ ; voir Variétés, lemme
\ref{varieties-lemma-numerical-polynomial-from-euler}.
Puisque $U$ est connexe et que les fonctions
$u \mapsto \chi(Y_u, \mathcal{N}_u^{\otimes n})$
sont localement constantes (voir Catégories dérivées des schémas, lemme
\ref{perfect-lemma-chi-locally-constant-geometric}), on obtient le même
polynôme de Hilbert en tout point de $U$. Posons
$\mathcal{M} = \text{pr}_1^*\mathcal{N}
\otimes_{\mathcal{O}_{Y \times_U X}} \text{pr}_2^*\mathcal{L}$
sur $Y \times_U X$. Étant donné
$(f, \varphi) \in \mathit{Isom}_{\Polarizedstack}(\upsilon, \chi)(T)$
pour un schéma $T$ sur $U$, on a pour tout $t \in T$
$$
\chi(Y_t, (\text{id} \times f)^*\mathcal{M}^{\otimes n}) =
\chi(Y_t,
\mathcal{N}_t^{\otimes n} \otimes_{\mathcal{O}_{Y_t}}
f_t^*\mathcal{L}_t^{\otimes n}) =
\chi(Y_t, \mathcal{N}_t^{\otimes 2n}) = P(2n)
$$
où, dans l'égalité du milieu, nous utilisons l'isomorphisme
$\varphi : f^*\mathcal{L}_T \to \mathcal{N}_T$.
Posant $P'(t) = P(2t)$, on voit que le morphisme
$$
\mathit{Isom}_{\Polarizedstack}(\upsilon, \chi)
\longrightarrow
\mathit{Isom}_U(Y, X)
$$
(voir plus haut) a son image contenue dans l'intersection
$$
\mathit{Isom}_U(Y, X) \cap \mathit{Mor}^{P', \mathcal{M}}_U(Y, X)
$$
Cette intersection est une intersection de sous-espaces ouverts de
$\mathit{Mor}_U(Y, X)$ (voir le lemme \ref{lemma-Isom-in-Mor} et la
remarque \ref{remark-Mor-numerical}).
Or $\mathit{Mor}^{P', \mathcal{M}}_U(Y, X)$ est un espace algébrique
noethérien, puisqu'il est de présentation finie sur $U$ d'après le lemme
\ref{lemma-Mor-qc-over-base}. L'intersection est donc elle aussi un espace
algébrique noethérien. Comme le morphisme
$$
\mathit{Isom}_{\Polarizedstack}(\upsilon, \chi)
\longrightarrow
\mathit{Isom}_U(Y, X) \cap \mathit{Mor}^{P', \mathcal{M}}_U(Y, X)
$$
est affine (voir plus haut), on conclut.
\end{proof}

\begin{lemma}
\label{lemma-polarized-qs-lfp}
Le morphisme $\Polarizedstack \to \Spec(\mathbf{Z})$ est quasi-séparé et
localement de présentation finie.
\end{lemma}

\begin{proof}
Pour vérifier que $\Polarizedstack \to \Spec(\mathbf{Z})$ est
quasi-séparé, il faut montrer que sa diagonale est quasi-compacte et
quasi-séparée. Cela résulte immédiatement du lemme
\ref{lemma-polarized-diagonal-separated-fp}.
Pour prouver que $\Polarizedstack \to \Spec(\mathbf{Z})$ est localement de
présentation finie, il suffit de montrer que $\Polarizedstack$ préserve les
limites ; voir Limites de champs, proposition
\ref{stacks-limits-proposition-characterize-locally-finite-presentation}.
C'est Quotients, lemme \ref{quot-lemma-polarized-limits}.
\end{proof}

\begin{lemma}
\label{lemma-bounded-polarized}
Soit $n \geq 1$ un entier et soit $P$ un polynôme numérique.
Soit
$$
T \subset |\Polarizedstack|
$$
un sous-ensemble possédant la propriété suivante : pour tout $\xi \in T$,
il existe un corps $k$ et un objet $(X, \mathcal{L})$ de
$\Polarizedstack$ sur $k$ représentant $\xi$, tels que
\begin{enumerate}
\item le polynôme de Hilbert de $\mathcal{L}$ sur $X$ soit $P$, et
\item il existe une immersion fermée $i : X \to \mathbf{P}^n_k$ telle que
$i^*\mathcal{O}_{\mathbf{P}^n}(1) \cong \mathcal{L}$.
\end{enumerate}
Alors $T$ est un espace topologique noethérien, en particulier
quasi-compact.
\end{lemma}

\begin{proof}
Remarquons que $|\Polarizedstack|$ est un espace topologique localement
noethérien ; voir Morphismes de champs, lemme
\ref{stacks-morphisms-lemma-Noetherian-topology}
(on utilise également le fait que $\Spec(\mathbf{Z})$ est noethérien et
que $\Polarizedstack$ est donc un champ algébrique localement noethérien
d'après le lemme \ref{lemma-polarized-qs-lfp} et Morphismes de champs,
lemme
\ref{stacks-morphisms-lemma-locally-finite-type-locally-noetherian}).
Tout sous-ensemble quasi-compact de $|\Polarizedstack|$ est donc un espace
topologique noethérien, et tout sous-ensemble d'un tel espace est lui aussi
noethérien ; voir Topologie, lemmes
\ref{topology-lemma-finite-union-Noetherian} et
\ref{topology-lemma-Noetherian}.
Il ne reste donc qu'à trouver un sous-ensemble quasi-compact contenant
$T$.

\medskip\noindent
D'après le lemme \ref{lemma-hilb-proper-over-base}, l'espace algébrique
$$
H =
\Hilbfunctor^{P, \mathcal{O}(1)}_{\mathbf{P}^n_\mathbf{Z}/\Spec(\mathbf{Z})}
$$
est propre sur $\Spec(\mathbf{Z})$. D'après Quotients, lemme
\ref{quot-lemma-extend-hilb-to-spaces}\footnote{Nous verrons plus loin
(insérer ici une référence future) que $H$ est un schéma ; l'emploi de ce
lemme et de Quotients, lemme
\ref{quot-lemma-extend-polarized-to-spaces}, n'est donc pas nécessaire.},
le morphisme identité de $H$ correspond à un sous-espace fermé
$$
Z \subset \mathbf{P}^n_H
$$
qui est propre, plat et de présentation finie sur $H$, tel que la
restriction $\mathcal{N} = \mathcal{O}(1)|_Z$ soit relativement ample sur
$Z/H$ et ait le polynôme de Hilbert $P$ sur les fibres de $Z \to H$.
En particulier, la paire $(Z \to H, \mathcal{N})$ définit un morphisme
$$
H \longrightarrow \Polarizedstack
$$
qui envoie un morphisme de schémas $U \to H$ sur le morphisme classifiant
de la famille $(Z_U \to U, \mathcal{N}_U)$ ; voir Quotients, lemme
\ref{quot-lemma-extend-polarized-to-spaces}.
Comme $H$ est un espace algébrique noethérien (puisqu'il est propre sur
$\mathbf{Z})$), on voit que $|H|$ est noethérien, donc quasi-compact.
L'application
$$
|H| \longrightarrow |\Polarizedstack|
$$
est continue ; son image est donc quasi-compacte.
Il suffit ainsi de prouver que $T$ est contenu dans l'image de
$|H| \to |\Polarizedstack|$.
Or les hypothèses (1) et (2) expriment exactement ce fait : tout choix
d'une immersion fermée $i : X \to \mathbf{P}^n_k$ telle que
$i^*\mathcal{O}_{\mathbf{P}^n}(1) \cong \mathcal{L}$ fournit un point à
valeurs dans $k$ de $H$, par l'interprétation modulaire de $H$.
Cela achève la démonstration du lemme.
\end{proof}






\section{Propriétés des modules de complexes sur un morphisme propre}
\label{section-complexes}

\noindent
Soit $f : X \to B$ un morphisme d'espaces algébriques propre, plat et de
présentation finie. Alors le champ $\Complexesstack_{X/B}$ qui paramètre
les complexes relativement parfaits dont les groupes d'auto-extensions de
degré négatif sont nuls est algébrique. Voir Quotients, théorème
\ref{quot-theorem-complexes-algebraic}.

\begin{lemma}
\label{lemma-complexes-diagonal-affine-fp}
La diagonale de $\Complexesstack_{X/B}$ sur $B$ est affine et de
présentation finie.
\end{lemma}

\begin{proof}
La représentabilité de la diagonale par des espaces algébriques a été
établie dans Quotients, lemme \ref{quot-lemma-complexes-diagonal}.
Il ressort de la démonstration qu'il faut montrer ceci : étant donnés un
schéma $T$ sur $B$ et des objets $E, E' \in D(\mathcal{O}_{X_T})$ tels que
$(T, E)$ et $(T, E')$ soient des objets de la catégorie fibre de
$\Complexesstack_{X/B}$ sur $T$, le morphisme
$\mathit{Isom}(E, E') \to T$ est affine et de présentation finie.
Ici, $\mathit{Isom}(E, E')$ est le foncteur
$$
(\Sch/T)^{opp} \to \textit{Sets},\quad
T' \mapsto \{\varphi : E_{T'} \to E'_{T'}
\text{ isomorphisme dans }D(\mathcal{O}_{X_{T'}})\}
$$
où $E_{T'}$ et $E'_{T'}$ sont les images inverses dérivées de $E$ et de
$E'$ sur $X_{T'}$. Considérons le foncteur
$H = \SheafHom(E, E')$ défini par la règle
$$
(\Sch/T)^{opp} \to \textit{Sets},\quad
T' \mapsto \Hom_{\mathcal{O}_{X_{T'}}}(E_T, E'_T)
$$
D'après Quotients, lemme
\ref{quot-lemma-complexes-open-neg-exts-vanishing}, c'est un espace
algébrique affine et de présentation finie sur $T$.
Il en va de même de $H' = \SheafHom(E', E)$,
$I = \SheafHom(E, E)$ et $I' = \SheafHom(E', E')$.
On voit donc que
$$
\mathit{Isom}(E, E') = (H' \times_T H) \times_{c, I \times_T I', \sigma} T
$$
où $c(\varphi', \varphi) =
(\varphi \circ \varphi', \varphi' \circ \varphi)$ et
$\sigma = (\text{id}, \text{id})$ (comparer à la démonstration de
Quotients, proposition \ref{quot-proposition-isom}).
Ainsi, $\mathit{Isom}(E, E')$ est affine sur $T$ comme produit fibré de
schémas affines sur $T$. De même, $\mathit{Isom}(E, E')$ est de
présentation finie sur $T$.
\end{proof}

\begin{lemma}
\label{lemma-complexes-qs-lfp}
Le morphisme $\Complexesstack_{X/B} \to B$ est quasi-séparé et localement
de présentation finie.
\end{lemma}

\begin{proof}
Pour vérifier que $\Complexesstack_{X/B} \to B$ est quasi-séparé, il faut
montrer que sa diagonale est quasi-compacte et quasi-séparée.
Cela résulte immédiatement du lemme
\ref{lemma-complexes-diagonal-affine-fp}.
Pour prouver que $\Complexesstack_{X/B} \to B$ est localement de
présentation finie, il faut montrer que $\Complexesstack_{X/B} \to B$
préserve les limites ; voir Limites de champs, proposition
\ref{stacks-limits-proposition-characterize-locally-finite-presentation}.
Cela résulte de Quotients, lemme \ref{quot-lemma-complexes-limits}
(un petit détail est omis).
\end{proof}






\begin{multicols}{2}[\section{Autres chapitres}]
\noindent
Préliminaires
\begin{enumerate}
\item \hyperref[introduction-section-phantom]{Introduction}
\item \hyperref[conventions-section-phantom]{Conventions}
\item \hyperref[sets-section-phantom]{Théorie des ensembles}
\item \hyperref[categories-section-phantom]{Catégories}
\item \hyperref[topology-section-phantom]{Topologie}
\item \hyperref[sheaves-section-phantom]{Faisceaux sur les espaces}
\item \hyperref[sites-section-phantom]{Sites et faisceaux}
\item \hyperref[stacks-section-phantom]{Champs}
\item \hyperref[fields-section-phantom]{Corps}
\item \hyperref[algebra-section-phantom]{Algèbre commutative}
\item \hyperref[brauer-section-phantom]{Groupes de Brauer}
\item \hyperref[homology-section-phantom]{Algèbre homologique}
\item \hyperref[derived-section-phantom]{Catégories dérivées}
\item \hyperref[simplicial-section-phantom]{Méthodes simpliciales}
\item \hyperref[more-algebra-section-phantom]{Compléments d'algèbre}
\item \hyperref[smoothing-section-phantom]{Lissage des morphismes d'anneaux}
\item \hyperref[modules-section-phantom]{Faisceaux de modules}
\item \hyperref[sites-modules-section-phantom]{Modules sur les sites}
\item \hyperref[injectives-section-phantom]{Objets injectifs}
\item \hyperref[cohomology-section-phantom]{Cohomologie des faisceaux}
\item \hyperref[sites-cohomology-section-phantom]{Cohomologie sur les sites}
\item \hyperref[dga-section-phantom]{Algèbre différentielle graduée}
\item \hyperref[dpa-section-phantom]{Algèbre à puissances divisées}
\item \hyperref[sdga-section-phantom]{Faisceaux différentiels gradués}
\item \hyperref[hypercovering-section-phantom]{Hyperrecouvrements}
\end{enumerate}
Schémas
\begin{enumerate}
\setcounter{enumi}{25}
\item \hyperref[schemes-section-phantom]{Schémas}
\item \hyperref[constructions-section-phantom]{Constructions de schémas}
\item \hyperref[properties-section-phantom]{Propriétés des schémas}
\item \hyperref[morphisms-section-phantom]{Morphismes de schémas}
\item \hyperref[coherent-section-phantom]{Cohomologie des schémas}
\item \hyperref[divisors-section-phantom]{Diviseurs}
\item \hyperref[limits-section-phantom]{Limites de schémas}
\item \hyperref[varieties-section-phantom]{Variétés}
\item \hyperref[topologies-section-phantom]{Topologies sur les schémas}
\item \hyperref[descent-section-phantom]{Descente}
\item \hyperref[perfect-section-phantom]{Catégories dérivées de schémas}
\item \hyperref[more-morphisms-section-phantom]{Compléments sur les morphismes}
\item \hyperref[flat-section-phantom]{Compléments sur la platitude}
\item \hyperref[groupoids-section-phantom]{Schémas en groupoïdes}
\item \hyperref[more-groupoids-section-phantom]{Compléments sur les schémas en groupoïdes}
\item \hyperref[etale-section-phantom]{Morphismes étales de schémas}
\end{enumerate}
Thèmes de théorie des schémas
\begin{enumerate}
\setcounter{enumi}{41}
\item \hyperref[chow-section-phantom]{Homologie de Chow}
\item \hyperref[intersection-section-phantom]{Théorie de l'intersection}
\item \hyperref[pic-section-phantom]{Schémas de Picard des courbes}
\item \hyperref[weil-section-phantom]{Théories cohomologiques de Weil}
\item \hyperref[adequate-section-phantom]{Modules adéquats}
\item \hyperref[dualizing-section-phantom]{Complexes dualisants}
\item \hyperref[duality-section-phantom]{Dualité pour les schémas}
\item \hyperref[discriminant-section-phantom]{Discriminants et différentes}
\item \hyperref[derham-section-phantom]{Cohomologie de de Rham}
\item \hyperref[local-cohomology-section-phantom]{Cohomologie locale}
\item \hyperref[algebraization-section-phantom]{Géométrie algébrique et formelle}
\item \hyperref[curves-section-phantom]{Courbes algébriques}
\item \hyperref[resolve-section-phantom]{Résolution des surfaces}
\item \hyperref[models-section-phantom]{Réduction semi-stable}
\item \hyperref[functors-section-phantom]{Foncteurs et morphismes}
\item \hyperref[equiv-section-phantom]{Catégories dérivées de variétés}
\item \hyperref[pione-section-phantom]{Groupes fondamentaux de schémas}
\item \hyperref[etale-cohomology-section-phantom]{Cohomologie étale}
\item \hyperref[crystalline-section-phantom]{Cohomologie cristalline}
\item \hyperref[proetale-section-phantom]{Cohomologie pro-étale}
\item \hyperref[relative-cycles-section-phantom]{Cycles relatifs}
\item \hyperref[more-etale-section-phantom]{Compléments de cohomologie étale}
\item \hyperref[trace-section-phantom]{La formule des traces}
\end{enumerate}
Espaces algébriques
\begin{enumerate}
\setcounter{enumi}{64}
\item \hyperref[spaces-section-phantom]{Espaces algébriques}
\item \hyperref[spaces-properties-section-phantom]{Propriétés des espaces algébriques}
\item \hyperref[spaces-morphisms-section-phantom]{Morphismes d'espaces algébriques}
\item \hyperref[decent-spaces-section-phantom]{Espaces algébriques décents}
\item \hyperref[spaces-cohomology-section-phantom]{Cohomologie des espaces algébriques}
\item \hyperref[spaces-limits-section-phantom]{Limites d'espaces algébriques}
\item \hyperref[spaces-divisors-section-phantom]{Diviseurs sur les espaces algébriques}
\item \hyperref[spaces-over-fields-section-phantom]{Espaces algébriques sur des corps}
\item \hyperref[spaces-topologies-section-phantom]{Topologies sur les espaces algébriques}
\item \hyperref[spaces-descent-section-phantom]{Descente et espaces algébriques}
\item \hyperref[spaces-perfect-section-phantom]{Catégories dérivées d'espaces}
\item \hyperref[spaces-more-morphisms-section-phantom]{Compléments sur les morphismes d'espaces}
\item \hyperref[spaces-flat-section-phantom]{Platitude sur les espaces algébriques}
\item \hyperref[spaces-groupoids-section-phantom]{Groupoïdes dans les espaces algébriques}
\item \hyperref[spaces-more-groupoids-section-phantom]{Compléments sur les groupoïdes dans les espaces}
\item \hyperref[bootstrap-section-phantom]{Amorçage}
\item \hyperref[spaces-pushouts-section-phantom]{Sommes amalgamées d'espaces algébriques}
\end{enumerate}
Thèmes de géométrie
\begin{enumerate}
\setcounter{enumi}{81}
\item \hyperref[spaces-chow-section-phantom]{Groupes de Chow des espaces}
\item \hyperref[groupoids-quotients-section-phantom]{Quotients de groupoïdes}
\item \hyperref[spaces-more-cohomology-section-phantom]{Compléments sur la cohomologie des espaces}
\item \hyperref[spaces-simplicial-section-phantom]{Espaces simpliciaux}
\item \hyperref[spaces-duality-section-phantom]{Dualité pour les espaces}
\item \hyperref[formal-spaces-section-phantom]{Espaces algébriques formels}
\item \hyperref[restricted-section-phantom]{Algébrisation des espaces formels}
\item \hyperref[spaces-resolve-section-phantom]{Résolution des surfaces revisitée}
\end{enumerate}
Théorie des déformations
\begin{enumerate}
\setcounter{enumi}{89}
\item \hyperref[formal-defos-section-phantom]{Théorie formelle des déformations}
\item \hyperref[defos-section-phantom]{Théorie des déformations}
\item \hyperref[cotangent-section-phantom]{Le complexe cotangent}
\item \hyperref[examples-defos-section-phantom]{Problèmes de déformation}
\end{enumerate}
Champs algébriques
\begin{enumerate}
\setcounter{enumi}{93}
\item \hyperref[algebraic-section-phantom]{Champs algébriques}
\item \hyperref[examples-stacks-section-phantom]{Exemples de champs}
\item \hyperref[stacks-sheaves-section-phantom]{Faisceaux sur les champs algébriques}
\item \hyperref[criteria-section-phantom]{Critères de représentabilité}
\item \hyperref[artin-section-phantom]{Axiomes d'Artin}
\item \hyperref[quot-section-phantom]{Espaces de Quot et de Hilbert}
\item \hyperref[stacks-properties-section-phantom]{Propriétés des champs algébriques}
\item \hyperref[stacks-morphisms-section-phantom]{Morphismes de champs algébriques}
\item \hyperref[stacks-limits-section-phantom]{Limites de champs algébriques}
\item \hyperref[stacks-cohomology-section-phantom]{Cohomologie des champs algébriques}
\item \hyperref[stacks-perfect-section-phantom]{Catégories dérivées de champs}
\item \hyperref[stacks-introduction-section-phantom]{Introduction aux champs algébriques}
\item \hyperref[stacks-more-morphisms-section-phantom]{Compléments sur les morphismes de champs}
\item \hyperref[stacks-geometry-section-phantom]{La géométrie des champs}
\end{enumerate}
Thèmes de théorie des espaces de modules
\begin{enumerate}
\setcounter{enumi}{107}
\item \hyperref[moduli-section-phantom]{Champs de modules}
\item \hyperref[moduli-curves-section-phantom]{Espaces de modules de courbes}
\end{enumerate}
Divers
\begin{enumerate}
\setcounter{enumi}{109}
\item \hyperref[examples-section-phantom]{Exemples}
\item \hyperref[exercises-section-phantom]{Exercices}
\item \hyperref[guide-section-phantom]{Guide bibliographique}
\item \hyperref[desirables-section-phantom]{Desiderata}
\item \hyperref[coding-section-phantom]{Style de programmation}
\item \hyperref[obsolete-section-phantom]{Obsolète}
\item \hyperref[fdl-section-phantom]{Licence de documentation libre GNU}
\item \hyperref[index-section-phantom]{Index généré automatiquement}
\end{enumerate}
\end{multicols}


\bibliography{my}
\bibliographystyle{amsalpha}

\end{document}



























